\chapter{命令行接口}

% === 源文件: cli/acp.md ===
\section{acp}


运行与 OpenClaw Gateway 网关通信的 ACP（Agent Client Protocol）桥接器。

此命令通过 stdio 使用 ACP 协议与 IDE 通信，并通过 WebSocket 将提示转发到 Gateway 网关。它将 ACP 会话映射到 Gateway 网关会话键。

\subsection{用法}


\begin{lstlisting}[language=bash]
openclaw acp

# Remote Gateway
openclaw acp --url wss://gateway-host:18789 --token <token>

# Attach to an existing session key
openclaw acp --session agent:main:main

# Attach by label (must already exist)
openclaw acp --session-label "support inbox"

# Reset the session key before the first prompt
openclaw acp --session agent:main:main --reset-session
\end{lstlisting}


\subsection{ACP 客户端（调试）}


使用内置 ACP 客户端在没有 IDE 的情况下检查桥接器的安装完整性。
它会启动 ACP 桥接器并让你交互式输入提示。

\begin{lstlisting}[language=bash]
openclaw acp client

# Point the spawned bridge at a remote Gateway
openclaw acp client --server-args --url wss://gateway-host:18789 --token <token>

# Override the server command (default: openclaw)
openclaw acp client --server "node" --server-args openclaw.mjs acp --url ws://127.0.0.1:19001
\end{lstlisting}


\subsection{如何使用}


当 IDE（或其他客户端）使用 Agent Client Protocol 并且你希望它驱动 OpenClaw Gateway 网关会话时，请使用 ACP。

\begin{enumerate}
  \item 确保 Gateway 网关正在运行（本地或远程）。
  \item 配置 Gateway 网关目标（配置或标志）。
  \item 将你的 IDE 配置为通过 stdio 运行 \texttt{openclaw acp}。
\end{enumerate}


示例配置（持久化）：

\begin{lstlisting}[language=bash]
openclaw config set gateway.remote.url wss://gateway-host:18789
openclaw config set gateway.remote.token <token>
\end{lstlisting}


示例直接运行（不写入配置）：

\begin{lstlisting}[language=bash]
openclaw acp --url wss://gateway-host:18789 --token <token>
\end{lstlisting}


\subsection{选择智能体}


ACP 不直接选择智能体。它通过 Gateway 网关会话键进行路由。

使用智能体作用域的会话键来定位特定智能体：

\begin{lstlisting}[language=bash]
openclaw acp --session agent:main:main
openclaw acp --session agent:design:main
openclaw acp --session agent:qa:bug-123
\end{lstlisting}


每个 ACP 会话映射到单个 Gateway 网关会话键。一个智能体可以有多个会话；除非你覆盖键或标签，否则 ACP 默认使用隔离的 \texttt{acp:} 会话。

\subsection{Zed 编辑器设置}


在 \texttt{\textasciitilde{}/.config/zed/settings.json} 中添加自定义 ACP 智能体（或使用 Zed 的设置界面）：

\begin{lstlisting}[language=json]
{
  "agent_servers": {
    "OpenClaw ACP": {
      "type": "custom",
      "command": "openclaw",
      "args": ["acp"],
      "env": {}
    }
  }
}
\end{lstlisting}


要定位特定的 Gateway 网关或智能体：

\begin{lstlisting}[language=json]
{
  "agent_servers": {
    "OpenClaw ACP": {
      "type": "custom",
      "command": "openclaw",
      "args": [
        "acp",
        "--url",
        "wss://gateway-host:18789",
        "--token",
        "<token>",
        "--session",
        "agent:design:main"
      ],
      "env": {}
    }
  }
}
\end{lstlisting}


在 Zed 中，打开 Agent 面板并选择"OpenClaw ACP"来开始一个会话。

\subsection{会话映射}


默认情况下，ACP 会话获得一个带有 \texttt{acp:} 前缀的隔离 Gateway 网关会话键。
要重用已知会话，请传递会话键或标签：

\begin{itemize}
  \item \texttt{--session }：使用特定的 Gateway 网关会话键。
  \item \texttt{--session-label }：通过标签解析现有会话。
  \item \texttt{--reset-session}：为该键生成新的会话 ID（相同键，新对话记录）。
\end{itemize}


如果你的 ACP 客户端支持元数据，你可以按会话覆盖：

\begin{lstlisting}[language=json]
{
  "_meta": {
    "sessionKey": "agent:main:main",
    "sessionLabel": "support inbox",
    "resetSession": true
  }
}
\end{lstlisting}


在 \href{/concepts/session}{/concepts/session} 了解更多关于会话键的信息。

\subsection{选项}


\begin{itemize}
  \item \texttt{--url }：Gateway 网关 WebSocket URL（配置后默认为 gateway.remote.url）。
  \item \texttt{--token }：Gateway 网关认证令牌。
  \item \texttt{--password }：Gateway 网关认证密码。
  \item \texttt{--session }：默认会话键。
  \item \texttt{--session-label }：要解析的默认会话标签。
  \item \texttt{--require-existing}：如果会话键/标签不存在则失败。
  \item \texttt{--reset-session}：在首次使用前重置会话键。
  \item \texttt{--no-prefix-cwd}：不在提示前添加工作目录前缀。
  \item \texttt{--verbose, -v}：向 stderr 输出详细日志。
\end{itemize}


\subsubsection{acp client 选项}


\begin{itemize}
  \item \texttt{--cwd }：ACP 会话的工作目录。
  \item \texttt{--server }：ACP 服务器命令（默认：\texttt{openclaw}）。
  \item \texttt{--server-args }：传递给 ACP 服务器的额外参数。
  \item \texttt{--server-verbose}：启用 ACP 服务器的详细日志。
  \item \texttt{--verbose, -v}：详细客户端日志。
\end{itemize}



\clearpage

% === 源文件: cli/agent.md ===
\section{openclaw agent}


通过 Gateway 网关运行智能体回合（使用 \texttt{--local} 进行嵌入式运行）。使用 \texttt{--agent } 直接指定已配置的智能体。

相关内容：

\begin{itemize}
  \item 智能体发送工具：\href{/tools/agent-send}{Agent send}
\end{itemize}


\subsection{示例}


\begin{lstlisting}[language=bash]
openclaw agent --to +15555550123 --message "status update" --deliver
openclaw agent --agent ops --message "Summarize logs"
openclaw agent --session-id 1234 --message "Summarize inbox" --thinking medium
openclaw agent --agent ops --message "Generate report" --deliver --reply-channel slack --reply-to "#reports"
\end{lstlisting}



\clearpage

% === 源文件: cli/agents.md ===
\section{openclaw agents}


管理隔离的智能体（工作区 + 认证 + 路由）。

相关内容：

\begin{itemize}
  \item 多智能体路由：\href{/concepts/multi-agent}{多智能体路由}
  \item 智能体工作区：\href{/concepts/agent-workspace}{智能体工作区}
\end{itemize}


\subsection{示例}


\begin{lstlisting}[language=bash]
openclaw agents list
openclaw agents add work --workspace ~/.openclaw/workspace-work
openclaw agents set-identity --workspace ~/.openclaw/workspace --from-identity
openclaw agents set-identity --agent main --avatar avatars/openclaw.png
openclaw agents delete work
\end{lstlisting}


\subsection{身份文件}


每个智能体工作区可以在工作区根目录包含一个 \texttt{IDENTITY.md}：

\begin{itemize}
  \item 示例路径：\texttt{\textasciitilde{}/.openclaw/workspace/IDENTITY.md}
  \item \texttt{set-identity --from-identity} 从工作区根目录读取（或从显式指定的 \texttt{--identity-file} 读取）
\end{itemize}


头像路径相对于工作区根目录解析。

\subsection{设置身份}


\texttt{set-identity} 将字段写入 \texttt{agents.list[].identity}：

\begin{itemize}
  \item \texttt{name}
  \item \texttt{theme}
  \item \texttt{emoji}
  \item \texttt{avatar}（工作区相对路径、http(s) URL 或 data URI）
\end{itemize}


从 \texttt{IDENTITY.md} 加载：

\begin{lstlisting}[language=bash]
openclaw agents set-identity --workspace ~/.openclaw/workspace --from-identity
\end{lstlisting}


显式覆盖字段：

\begin{lstlisting}[language=bash]
openclaw agents set-identity --agent main --name "OpenClaw" --emoji "🦞" --avatar avatars/openclaw.png
\end{lstlisting}


配置示例：

\begin{lstlisting}[language=json]
{
  agents: {
    list: [
      {
        id: "main",
        identity: {
          name: "OpenClaw",
          theme: "space lobster",
          emoji: "🦞",
          avatar: "avatars/openclaw.png",
        },
      },
    ],
  },
}
\end{lstlisting}



\clearpage

% === 源文件: cli/approvals.md ===
\section{openclaw approvals}


管理\textbf{本地主机}、\textbf{Gateway 网关主机}或\textbf{节点主机}的执行审批。
默认情况下，命令针对磁盘上的本地审批文件。使用 \texttt{--gateway} 可针对 Gateway 网关，使用 \texttt{--node} 可针对特定节点。

相关内容：

\begin{itemize}
  \item 执行审批：\href{/tools/exec-approvals}{执行审批}
  \item 节点：\href{/nodes}{节点}
\end{itemize}


\subsection{常用命令}


\begin{lstlisting}[language=bash]
openclaw approvals get
openclaw approvals get --node <id|name|ip>
openclaw approvals get --gateway
\end{lstlisting}


\subsection{从文件替换审批}


\begin{lstlisting}[language=bash]
openclaw approvals set --file ./exec-approvals.json
openclaw approvals set --node <id|name|ip> --file ./exec-approvals.json
openclaw approvals set --gateway --file ./exec-approvals.json
\end{lstlisting}


\subsection{允许列表辅助命令}


\begin{lstlisting}[language=bash]
openclaw approvals allowlist add "~/Projects/**/bin/rg"
openclaw approvals allowlist add --agent main --node <id|name|ip> "/usr/bin/uptime"
openclaw approvals allowlist add --agent "*" "/usr/bin/uname"

openclaw approvals allowlist remove "~/Projects/**/bin/rg"
\end{lstlisting}


\subsection{注意事项}


\begin{itemize}
  \item \texttt{--node} 使用与 \texttt{openclaw nodes} 相同的解析器（id、name、ip 或 id 前缀）。
  \item \texttt{--agent} 默认为 \texttt{"*"}，表示适用于所有智能体。
  \item 节点主机必须公开 \texttt{system.execApprovals.get/set}（macOS 应用或无头节点主机）。
  \item 审批文件按主机存储在 \texttt{\textasciitilde{}/.openclaw/exec-approvals.json}。
\end{itemize}



\clearpage

% === 源文件: cli/browser.md ===
\section{openclaw browser}


管理 OpenClaw 的浏览器控制服务器并运行浏览器操作（标签页、快照、截图、导航、点击、输入）。

相关：

\begin{itemize}
  \item 浏览器工具 + API：\href{/tools/browser}{浏览器工具}
  \item Chrome 扩展中继：\href{/tools/chrome-extension}{Chrome 扩展}
\end{itemize}


\subsection{通用标志}


\begin{itemize}
  \item \texttt{--url }：Gateway 网关 WebSocket URL（默认从配置获取）。
  \item \texttt{--token }：Gateway 网关令牌（如果需要）。
  \item \texttt{--timeout }：请求超时（毫秒）。
  \item \texttt{--browser-profile }：选择浏览器配置文件（默认从配置获取）。
  \item \texttt{--json}：机器可读输出（在支持的地方）。
\end{itemize}


\subsection{快速开始（本地）}


\begin{lstlisting}[language=bash]
openclaw browser --browser-profile chrome tabs
openclaw browser --browser-profile openclaw start
openclaw browser --browser-profile openclaw open https://example.com
openclaw browser --browser-profile openclaw snapshot
\end{lstlisting}


\subsection{配置文件}


配置文件是命名的浏览器路由配置。实际上：

\begin{itemize}
  \item \texttt{openclaw}：启动/附加到专用的 OpenClaw 管理的 Chrome 实例（隔离的用户数据目录）。
  \item \texttt{chrome}：通过 Chrome 扩展中继控制你现有的 Chrome 标签页。
\end{itemize}


\begin{lstlisting}[language=bash]
openclaw browser profiles
openclaw browser create-profile --name work --color "#FF5A36"
openclaw browser delete-profile --name work
\end{lstlisting}


使用特定配置文件：

\begin{lstlisting}[language=bash]
openclaw browser --browser-profile work tabs
\end{lstlisting}


\subsection{标签页}


\begin{lstlisting}[language=bash]
openclaw browser tabs
openclaw browser open https://docs.openclaw.ai
openclaw browser focus <targetId>
openclaw browser close <targetId>
\end{lstlisting}


\subsection{快照 / 截图 / 操作}


快照：

\begin{lstlisting}[language=bash]
openclaw browser snapshot
\end{lstlisting}


截图：

\begin{lstlisting}[language=bash]
openclaw browser screenshot
\end{lstlisting}


导航/点击/输入（基于 ref 的 UI 自动化）：

\begin{lstlisting}[language=bash]
openclaw browser navigate https://example.com
openclaw browser click <ref>
openclaw browser type <ref> "hello"
\end{lstlisting}


\subsection{Chrome 扩展中继（通过工具栏按钮附加）}


此模式让智能体控制你手动附加的现有 Chrome 标签页（不会自动附加）。

将未打包的扩展安装到稳定路径：

\begin{lstlisting}[language=bash]
openclaw browser extension install
openclaw browser extension path
\end{lstlisting}


然后 Chrome → \texttt{chrome://extensions} → 启用"开发者模式" → "加载已解压的扩展程序" → 选择打印的文件夹。

完整指南：\href{/tools/chrome-extension}{Chrome 扩展}

\subsection{远程浏览器控制（node host 代理）}


如果 Gateway 网关与浏览器运行在不同的机器上，在有 Chrome/Brave/Edge/Chromium 的机器上运行 \textbf{node host}。Gateway 网关会将浏览器操作代理到该节点（无需单独的浏览器控制服务器）。

使用 \texttt{gateway.nodes.browser.mode} 控制自动路由，使用 \texttt{gateway.nodes.browser.node} 在连接多个节点时固定特定节点。

安全 + 远程设置：\href{/tools/browser}{浏览器工具}、\href{/gateway/remote}{远程访问}、\href{/gateway/tailscale}{Tailscale}、\href{/gateway/security}{安全}


\clearpage

% === 源文件: cli/channels.md ===
\section{openclaw channels}


管理 Gateway 网关上的聊天渠道账户及其运行时状态。

相关文档：

\begin{itemize}
  \item 渠道指南：\href{/channels/index}{渠道}
  \item Gateway 网关配置：\href{/gateway/configuration}{配置}
\end{itemize}


\subsection{常用命令}


\begin{lstlisting}[language=bash]
openclaw channels list
openclaw channels status
openclaw channels capabilities
openclaw channels capabilities --channel discord --target channel:123
openclaw channels resolve --channel slack "#general" "@jane"
openclaw channels logs --channel all
\end{lstlisting}


\subsection{添加/删除账户}


\begin{lstlisting}[language=bash]
openclaw channels add --channel telegram --token <bot-token>
openclaw channels remove --channel telegram --delete
\end{lstlisting}


提示：\texttt{openclaw channels add --help} 显示每个渠道的标志（token、app token、signal-cli 路径等）。

\subsection{登录/登出（交互式）}


\begin{lstlisting}[language=bash]
openclaw channels login --channel whatsapp
openclaw channels logout --channel whatsapp
\end{lstlisting}


\subsection{故障排除}


\begin{itemize}
  \item 运行 \texttt{openclaw status --deep} 进行全面探测。
  \item 使用 \texttt{openclaw doctor} 获取引导式修复。
  \item \texttt{openclaw channels list} 输出 \texttt{Claude: HTTP 403 ... user:profile} → 用量快照需要 \texttt{user:profile} 权限范围。使用 \texttt{--no-usage}，或提供 claude.ai 会话密钥（\texttt{CLAUDE\_WEB\_SESSION\_KEY} / \texttt{CLAUDE\_WEB\_COOKIE}），或通过 Claude Code CLI 重新授权。
\end{itemize}


\subsection{能力探测}


获取提供商能力提示（可用的 intents/scopes）以及静态功能支持：

\begin{lstlisting}[language=bash]
openclaw channels capabilities
openclaw channels capabilities --channel discord --target channel:123
\end{lstlisting}


说明：

\begin{itemize}
  \item \texttt{--channel} 是可选的；省略它可列出所有渠道（包括扩展）。
  \item \texttt{--target} 接受 \texttt{channel:} 或原始数字频道 id，仅适用于 Discord。
  \item 探测是特定于提供商的：Discord intents + 可选的频道权限；Slack bot + user scopes；Telegram bot 标志 + webhook；Signal daemon 版本；MS Teams app token + Graph roles/scopes（在已知处标注）。没有探测功能的渠道报告 \texttt{Probe: unavailable}。
\end{itemize}


\subsection{解析名称为 ID}


使用提供商目录将渠道/用户名称解析为 ID：

\begin{lstlisting}[language=bash]
openclaw channels resolve --channel slack "#general" "@jane"
openclaw channels resolve --channel discord "My Server/#support" "@someone"
openclaw channels resolve --channel matrix "Project Room"
\end{lstlisting}


说明：

\begin{itemize}
  \item 使用 \texttt{--kind user|group|auto} 强制指定目标类型。
  \item 当多个条目共享相同名称时，解析优先选择活跃的匹配项。
\end{itemize}



\clearpage

% === 源文件: cli/config.md ===
\section{openclaw config}


配置辅助命令：通过路径获取/设置/取消设置值。不带子命令运行将打开
配置向导（与 \texttt{openclaw configure} 相同）。

\subsection{示例}


\begin{lstlisting}[language=bash]
openclaw config get browser.executablePath
openclaw config set browser.executablePath "/usr/bin/google-chrome"
openclaw config set agents.defaults.heartbeat.every "2h"
openclaw config set agents.list[0].tools.exec.node "node-id-or-name"
openclaw config unset tools.web.search.apiKey
\end{lstlisting}


\subsection{路径}


路径使用点号或括号表示法：

\begin{lstlisting}[language=bash]
openclaw config get agents.defaults.workspace
openclaw config get agents.list[0].id
\end{lstlisting}


使用智能体列表索引来定位特定智能体：

\begin{lstlisting}[language=bash]
openclaw config get agents.list
openclaw config set agents.list[1].tools.exec.node "node-id-or-name"
\end{lstlisting}


\subsection{值}


值会尽可能解析为 JSON5；否则将被视为字符串。
使用 \texttt{--json} 强制要求 JSON5 解析。

\begin{lstlisting}[language=bash]
openclaw config set agents.defaults.heartbeat.every "0m"
openclaw config set gateway.port 19001 --json
openclaw config set channels.whatsapp.groups '["*"]' --json
\end{lstlisting}


编辑后请重启 Gateway 网关。


\clearpage

% === 源文件: cli/configure.md ===
\section{openclaw configure}


用于设置凭证、设备和智能体默认值的交互式提示。

注意：\textbf{模型}部分现在包含一个用于 \texttt{agents.defaults.models} 允许列表的多选项（显示在 \texttt{/model} 和模型选择器中的内容）。

提示：不带子命令的 \texttt{openclaw config} 会打开相同的向导。使用 \texttt{openclaw config get|set|unset} 进行非交互式编辑。

相关内容：

\begin{itemize}
  \item Gateway 网关配置参考：\href{/gateway/configuration}{配置}
  \item Config CLI：\href{/cli/config}{Config}
\end{itemize}


注意事项：

\begin{itemize}
  \item 选择 Gateway 网关运行位置始终会更新 \texttt{gateway.mode}。如果这是你唯一需要的，可以不选择其他部分直接选择"继续"。
  \item 面向渠道的服务（Slack/Discord/Matrix/Microsoft Teams）在设置期间会提示输入频道/房间允许列表。你可以输入名称或 ID；向导会尽可能将名称解析为 ID。
\end{itemize}


\subsection{示例}


\begin{lstlisting}[language=bash]
openclaw configure
openclaw configure --section models --section channels
\end{lstlisting}



\clearpage

% === 源文件: cli/cron.md ===
\section{openclaw cron}


管理 Gateway 网关调度器的 cron 作业。

相关内容：

\begin{itemize}
  \item Cron 作业：\href{/automation/cron-jobs}{Cron 作业}
\end{itemize}


提示：运行 \texttt{openclaw cron --help} 查看完整的命令集。

说明：隔离式 \texttt{cron add} 任务默认使用 \texttt{--announce} 投递摘要。使用 \texttt{--no-deliver} 仅内部运行。
\texttt{--deliver} 仍作为 \texttt{--announce} 的弃用别名保留。

说明：一次性（\texttt{--at}）任务成功后默认删除。使用 \texttt{--keep-after-run} 保留。

\subsection{常见编辑}


更新投递设置而不更改消息：

\begin{lstlisting}[language=bash]
openclaw cron edit <job-id> --announce --channel telegram --to "123456789"
\end{lstlisting}


为隔离的作业禁用投递：

\begin{lstlisting}[language=bash]
openclaw cron edit <job-id> --no-deliver
\end{lstlisting}



\clearpage

% === 源文件: cli/dashboard.md ===
\section{openclaw dashboard}


使用当前认证信息打开控制界面。

\begin{lstlisting}[language=bash]
openclaw dashboard
openclaw dashboard --no-open
\end{lstlisting}



\clearpage

% === 源文件: cli/devices.md ===
\section{openclaw devices}


管理设备配对请求和设备范围的 token。

\subsection{命令}


\subsubsection{openclaw devices list}


列出待处理的配对请求和已配对的设备。

\begin{lstlisting}
openclaw devices list
openclaw devices list --json
\end{lstlisting}


\subsubsection{openclaw devices approve }


批准待处理的设备配对请求。

\begin{lstlisting}
openclaw devices approve <requestId>
\end{lstlisting}


\subsubsection{openclaw devices reject }


拒绝待处理的设备配对请求。

\begin{lstlisting}
openclaw devices reject <requestId>
\end{lstlisting}


\subsubsection{openclaw devices rotate --device  --role  [--scope ]}


为特定角色轮换设备 token（可选更新 scope）。

\begin{lstlisting}
openclaw devices rotate --device <deviceId> --role operator --scope operator.read --scope operator.write
\end{lstlisting}


\subsubsection{openclaw devices revoke --device  --role }


为特定角色撤销设备 token。

\begin{lstlisting}
openclaw devices revoke --device <deviceId> --role node
\end{lstlisting}


\subsection{通用选项}


\begin{itemize}
  \item \texttt{--url }：Gateway 网关 WebSocket URL（配置后默认使用 \texttt{gateway.remote.url}）。
  \item \texttt{--token }：Gateway 网关 token（如需要）。
  \item \texttt{--password }：Gateway 网关密码（密码认证）。
  \item \texttt{--timeout }：RPC 超时。
  \item \texttt{--json}：JSON 输出（推荐用于脚本）。
\end{itemize}


\subsection{注意事项}


\begin{itemize}
  \item Token 轮换会返回新 token（敏感信息）。请像对待密钥一样对待它。
  \item 这些命令需要 \texttt{operator.pairing}（或 \texttt{operator.admin}）scope。
\end{itemize}



\clearpage

% === 源文件: cli/directory.md ===
\section{openclaw directory}


对支持目录功能的渠道进行查找（联系人/对等方、群组和"我"）。

\subsection{通用参数}


\begin{itemize}
  \item \texttt{--channel }：渠道 ID/别名（配置了多个渠道时为必填；仅配置一个渠道时自动选择）
  \item \texttt{--account }：账号 ID（默认：渠道默认账号）
  \item \texttt{--json}：输出 JSON 格式
\end{itemize}


\subsection{说明}


\begin{itemize}
  \item \texttt{directory} 用于帮助你查找可粘贴到其他命令中的 ID（特别是 \texttt{openclaw message send --target ...}）。
  \item 对于许多渠道，结果来源于配置（允许列表/已配置的群组），而非实时的提供商目录。
  \item 默认输出为以制表符分隔的 \texttt{id}（有时包含 \texttt{name}）；脚本中请使用 \texttt{--json}。
\end{itemize}


\subsection{将结果用于 message send}


\begin{lstlisting}[language=bash]
openclaw directory peers list --channel slack --query "U0"
openclaw message send --channel slack --target user:U012ABCDEF --message "hello"
\end{lstlisting}


\subsection{ID 格式（按渠道）}


\begin{itemize}
  \item WhatsApp：\texttt{+15551234567}（私聊），\texttt{1234567890-1234567890@g.us}（群组）
  \item Telegram：\texttt{@username} 或数字聊天 ID；群组为数字 ID
  \item Slack：\texttt{user:U…} 和 \texttt{channel:C…}
  \item Discord：\texttt{user:} 和 \texttt{channel:}
  \item Matrix（插件）：\texttt{user:@user:server}、\texttt{room:!roomId:server} 或 \texttt{\#alias:server}
  \item Microsoft Teams（插件）：\texttt{user:} 和 \texttt{conversation:}
  \item Zalo（插件）：用户 ID（Bot API）
  \item Zalo Personal / \texttt{zalouser}（插件）：来自 \texttt{zca} 的会话 ID（私聊/群组）（\texttt{me}、\texttt{friend list}、\texttt{group list}）
\end{itemize}


\subsection{Self（"我"）}


\begin{lstlisting}[language=bash]
openclaw directory self --channel zalouser
\end{lstlisting}


\subsection{Peers（联系人/用户）}


\begin{lstlisting}[language=bash]
openclaw directory peers list --channel zalouser
openclaw directory peers list --channel zalouser --query "name"
openclaw directory peers list --channel zalouser --limit 50
\end{lstlisting}


\subsection{群组}


\begin{lstlisting}[language=bash]
openclaw directory groups list --channel zalouser
openclaw directory groups list --channel zalouser --query "work"
openclaw directory groups members --channel zalouser --group-id <id>
\end{lstlisting}



\clearpage

% === 源文件: cli/dns.md ===
\section{openclaw dns}


用于广域设备发现（Tailscale + CoreDNS）的 DNS 辅助工具。目前专注于 macOS + Homebrew CoreDNS。

相关内容：

\begin{itemize}
  \item Gateway 网关设备发现：\href{/gateway/discovery}{设备发现}
  \item 广域设备发现配置：\href{/gateway/configuration}{配置}
\end{itemize}


\subsection{设置}


\begin{lstlisting}[language=bash]
openclaw dns setup
openclaw dns setup --apply
\end{lstlisting}



\clearpage

% === 源文件: cli/docs.md ===
\section{openclaw docs}


搜索实时文档索引。

\begin{lstlisting}[language=bash]
openclaw docs browser extension
openclaw docs sandbox allowHostControl
\end{lstlisting}



\clearpage

% === 源文件: cli/doctor.md ===
\section{openclaw doctor}


Gateway 网关和渠道的健康检查 + 快速修复。

相关内容：

\begin{itemize}
  \item 故障排除：\href{/gateway/troubleshooting}{故障排除}
  \item 安全审计：\href{/gateway/security}{安全}
\end{itemize}


\subsection{示例}


\begin{lstlisting}[language=bash]
openclaw doctor
openclaw doctor --repair
openclaw doctor --deep
\end{lstlisting}


注意事项：

\begin{itemize}
  \item 交互式提示（如钥匙串/OAuth 修复）仅在 stdin 是 TTY 且\textbf{未}设置 \texttt{--non-interactive} 时运行。无头运行（cron、Telegram、无终端）将跳过提示。
  \item \texttt{--fix}（\texttt{--repair} 的别名）会将备份写入 \texttt{\textasciitilde{}/.openclaw/openclaw.json.bak}，并删除未知的配置键，同时列出每个删除项。
\end{itemize}


\subsection{macOS：launchctl 环境变量覆盖}


如果你之前运行过 \texttt{launchctl setenv OPENCLAW\_GATEWAY\_TOKEN ...}（或 \texttt{...PASSWORD}），该值会覆盖你的配置文件，并可能导致持续的"未授权"错误。

\begin{lstlisting}[language=bash]
launchctl getenv OPENCLAW_GATEWAY_TOKEN
launchctl getenv OPENCLAW_GATEWAY_PASSWORD

launchctl unsetenv OPENCLAW_GATEWAY_TOKEN
launchctl unsetenv OPENCLAW_GATEWAY_PASSWORD
\end{lstlisting}



\clearpage

% === 源文件: cli/gateway.md ===
\section{Gateway 网关 CLI}


Gateway 网关是 OpenClaw 的 WebSocket 服务器（渠道、节点、会话、hooks）。

本页中的子命令位于 \texttt{openclaw gateway …} 下。

相关文档：

\begin{itemize}
  \item \href{/gateway/bonjour}{/gateway/bonjour}
  \item \href{/gateway/discovery}{/gateway/discovery}
  \item \href{/gateway/configuration}{/gateway/configuration}
\end{itemize}


\subsection{运行 Gateway 网关}


运行本地 Gateway 网关进程：

\begin{lstlisting}[language=bash]
openclaw gateway
\end{lstlisting}


前台运行别名：

\begin{lstlisting}[language=bash]
openclaw gateway run
\end{lstlisting}


注意事项：

\begin{itemize}
  \item 默认情况下，除非在 \texttt{\textasciitilde{}/.openclaw/openclaw.json} 中设置了 \texttt{gateway.mode=local}，否则 Gateway 网关将拒绝启动。使用 \texttt{--allow-unconfigured} 进行临时/开发运行。
  \item 在没有认证的情况下绑定到 loopback 之外的地址会被阻止（安全护栏）。
  \item \texttt{SIGUSR1} 在授权时触发进程内重启（启用 \texttt{commands.restart} 或使用 gateway 工具/config apply/update）。
  \item \texttt{SIGINT}/\texttt{SIGTERM} 处理程序会停止 Gateway 网关进程，但不会恢复任何自定义终端状态。如果你用 TUI 或 raw-mode 输入包装 CLI，请在退出前恢复终端。
\end{itemize}


\subsubsection{选项}


\begin{itemize}
  \item \texttt{--port }：WebSocket 端口（默认来自配置/环境变量；通常为 \texttt{18789}）。
  \item \texttt{--bind }：监听器绑定模式。
  \item \texttt{--auth }：认证模式覆盖。
  \item \texttt{--token }：令牌覆盖（同时为进程设置 \texttt{OPENCLAW\_GATEWAY\_TOKEN}）。
  \item \texttt{--password }：密码覆盖（同时为进程设置 \texttt{OPENCLAW\_GATEWAY\_PASSWORD}）。
  \item \texttt{--tailscale }：通过 Tailscale 暴露 Gateway 网关。
  \item \texttt{--tailscale-reset-on-exit}：关闭时重置 Tailscale serve/funnel 配置。
  \item \texttt{--allow-unconfigured}：允许在配置中没有 \texttt{gateway.mode=local} 的情况下启动 Gateway 网关。
  \item \texttt{--dev}：如果缺失则创建开发配置 + 工作区（跳过 BOOTSTRAP.md）。
  \item \texttt{--reset}：重置开发配置 + 凭证 + 会话 + 工作区（需要 \texttt{--dev}）。
  \item \texttt{--force}：启动前杀死所选端口上的任何现有监听器。
  \item \texttt{--verbose}：详细日志。
  \item \texttt{--claude-cli-logs}：仅在控制台显示 claude-cli 日志（并启用其 stdout/stderr）。
  \item \texttt{--ws-log }：WebSocket 日志样式（默认 \texttt{auto}）。
  \item \texttt{--compact}：\texttt{--ws-log compact} 的别名。
  \item \texttt{--raw-stream}：将原始模型流事件记录到 jsonl。
  \item \texttt{--raw-stream-path }：原始流 jsonl 路径。
\end{itemize}


\subsection{查询运行中的 Gateway 网关}


所有查询命令使用 WebSocket RPC。

输出模式：

\begin{itemize}
  \item 默认：人类可读（TTY 中带颜色）。
  \item \texttt{--json}：机器可读 JSON（无样式/进度指示器）。
  \item \texttt{--no-color}（或 \texttt{NO\_COLOR=1}）：禁用 ANSI 但保持人类可读布局。
\end{itemize}


共享选项（在支持的地方）：

\begin{itemize}
  \item \texttt{--url }：Gateway 网关 WebSocket URL。
  \item \texttt{--token }：Gateway 网关令牌。
  \item \texttt{--password }：Gateway 网关密码。
  \item \texttt{--timeout }：超时/预算（因命令而异）。
  \item \texttt{--expect-final}：等待"最终"响应（智能体调用）。
\end{itemize}


\subsubsection{gateway health}


\begin{lstlisting}[language=bash]
openclaw gateway health --url ws://127.0.0.1:18789
\end{lstlisting}


\subsubsection{gateway status}


\texttt{gateway status} 显示 Gateway 网关服务（launchd/systemd/schtasks）以及可选的 RPC 探测。

\begin{lstlisting}[language=bash]
openclaw gateway status
openclaw gateway status --json
\end{lstlisting}


选项：

\begin{itemize}
  \item \texttt{--url }：覆盖探测 URL。
  \item \texttt{--token }：探测的令牌认证。
  \item \texttt{--password }：探测的密码认证。
  \item \texttt{--timeout }：探测超时（默认 \texttt{10000}）。
  \item \texttt{--no-probe}：跳过 RPC 探测（仅服务视图）。
  \item \texttt{--deep}：也扫描系统级服务。
\end{itemize}


\subsubsection{gateway probe}


\texttt{gateway probe} 是"调试一切"命令。它始终探测：

\begin{itemize}
  \item 你配置的远程 Gateway 网关（如果设置了），以及
  \item localhost（loopback）\textbf{即使配置了远程也会探测}。
\end{itemize}


如果多个 Gateway 网关可达，它会打印所有。当你使用隔离的配置文件/端口（例如救援机器人）时支持多个 Gateway 网关，但大多数安装仍然运行单个 Gateway 网关。

\begin{lstlisting}[language=bash]
openclaw gateway probe
openclaw gateway probe --json
\end{lstlisting}


\paragraph{通过 SSH 远程（Mac 应用对等）}


macOS 应用的"通过 SSH 远程"模式使用本地端口转发，因此远程 Gateway 网关（可能仅绑定到 loopback）变得可以通过 \texttt{ws://127.0.0.1:} 访问。

CLI 等效命令：

\begin{lstlisting}[language=bash]
openclaw gateway probe --ssh user@gateway-host
\end{lstlisting}


选项：

\begin{itemize}
  \item \texttt{--ssh }：\texttt{user@host} 或 \texttt{user@host:port}（端口默认为 \texttt{22}）。
  \item \texttt{--ssh-identity }：身份文件。
  \item \texttt{--ssh-auto}：选择第一个发现的 Gateway 网关主机作为 SSH 目标（仅限局域网/WAB）。
\end{itemize}


配置（可选，用作默认值）：

\begin{itemize}
  \item \texttt{gateway.remote.sshTarget}
  \item \texttt{gateway.remote.sshIdentity}
\end{itemize}


\subsubsection{gateway call }


低级 RPC 辅助工具。

\begin{lstlisting}[language=bash]
openclaw gateway call status
openclaw gateway call logs.tail --params '{"sinceMs": 60000}'
\end{lstlisting}


\subsection{管理 Gateway 网关服务}


\begin{lstlisting}[language=bash]
openclaw gateway install
openclaw gateway start
openclaw gateway stop
openclaw gateway restart
openclaw gateway uninstall
\end{lstlisting}


注意事项：

\begin{itemize}
  \item \texttt{gateway install} 支持 \texttt{--port}、\texttt{--runtime}、\texttt{--token}、\texttt{--force}、\texttt{--json}。
  \item 生命周期命令接受 \texttt{--json} 用于脚本。
\end{itemize}


\subsection{发现 Gateway 网关（Bonjour）}


\texttt{gateway discover} 扫描 Gateway 网关信标（\texttt{\_openclaw-gw.\_tcp}）。

\begin{itemize}
  \item 组播 DNS-SD：\texttt{local.}
  \item 单播 DNS-SD（广域 Bonjour）：选择一个域（示例：\texttt{openclaw.internal.}）并设置分割 DNS + DNS 服务器；参见 \href{/gateway/bonjour}{/gateway/bonjour}
\end{itemize}


只有启用了 Bonjour 发现（默认）的 Gateway 网关才会广播信标。

广域发现记录包括（TXT）：

\begin{itemize}
  \item \texttt{role}（Gateway 网关角色提示）
  \item \texttt{transport}（传输提示，例如 \texttt{gateway}）
  \item \texttt{gatewayPort}（WebSocket 端口，通常为 \texttt{18789}）
  \item \texttt{sshPort}（SSH 端口；如果不存在则默认为 \texttt{22}）
  \item \texttt{tailnetDns}（MagicDNS 主机名，如果可用）
  \item \texttt{gatewayTls} / \texttt{gatewayTlsSha256}（TLS 启用 + 证书指纹）
  \item \texttt{cliPath}（远程安装的可选提示）
\end{itemize}


\subsubsection{gateway discover}


\begin{lstlisting}[language=bash]
openclaw gateway discover
\end{lstlisting}


选项：

\begin{itemize}
  \item \texttt{--timeout }：每个命令的超时（浏览/解析）；默认 \texttt{2000}。
  \item \texttt{--json}：机器可读输出（同时禁用样式/进度指示器）。
\end{itemize}


示例：

\begin{lstlisting}[language=bash]
openclaw gateway discover --timeout 4000
openclaw gateway discover --json | jq '.beacons[].wsUrl'
\end{lstlisting}



\clearpage

% === 源文件: cli/health.md ===
\section{openclaw health}


从运行中的 Gateway 网关获取健康状态。

\begin{lstlisting}[language=bash]
openclaw health
openclaw health --json
openclaw health --verbose
\end{lstlisting}


注意：

\begin{itemize}
  \item \texttt{--verbose} 运行实时探测，并在配置了多个账户时打印每个账户的耗时。
  \item 当配置了多个智能体时，输出包括每个智能体的会话存储。
\end{itemize}



\clearpage

% === 源文件: cli/hooks.md ===
\section{openclaw hooks}


管理智能体钩子（针对 \texttt{/new}、\texttt{/reset} 等命令以及 Gateway 网关启动的事件驱动自动化）。

相关内容：

\begin{itemize}
  \item 钩子：\href{/automation/hooks}{钩子}
  \item 插件钩子：\href{/tools/plugin\#plugin-hooks}{插件}
\end{itemize}


\subsection{列出所有钩子}


\begin{lstlisting}[language=bash]
openclaw hooks list
\end{lstlisting}


列出从工作区、托管目录和内置目录中发现的所有钩子。

\textbf{选项：}

\begin{itemize}
  \item \texttt{--eligible}：仅显示符合条件的钩子（满足要求）
  \item \texttt{--json}：以 JSON 格式输出
  \item \texttt{-v, --verbose}：显示详细信息，包括缺失的要求
\end{itemize}


\textbf{示例输出：}

\begin{lstlisting}
Hooks (3/3 ready)

Ready:
  🚀 boot-md ✓ - Run BOOT.md on gateway startup
  📝 command-logger ✓ - Log all command events to a centralized audit file
  💾 session-memory ✓ - Save session context to memory when /new command is issued
\end{lstlisting}


\textbf{示例（详细模式）：}

\begin{lstlisting}[language=bash]
openclaw hooks list --verbose
\end{lstlisting}


显示不符合条件的钩子缺失的要求。

\textbf{示例（JSON）：}

\begin{lstlisting}[language=bash]
openclaw hooks list --json
\end{lstlisting}


返回结构化 JSON，供程序化使用。

\subsection{获取钩子信息}


\begin{lstlisting}[language=bash]
openclaw hooks info <name>
\end{lstlisting}


显示特定钩子的详细信息。

\textbf{参数：}

\begin{itemize}
  \item `\texttt{：钩子名称（例如 }session-memory`）
\end{itemize}


\textbf{选项：}

\begin{itemize}
  \item \texttt{--json}：以 JSON 格式输出
\end{itemize}


\textbf{示例：}

\begin{lstlisting}[language=bash]
openclaw hooks info session-memory
\end{lstlisting}


\textbf{输出：}

\begin{lstlisting}
💾 session-memory ✓ Ready

Save session context to memory when /new command is issued

Details:
  Source: openclaw-bundled
  Path: /path/to/openclaw/hooks/bundled/session-memory/HOOK.md
  Handler: /path/to/openclaw/hooks/bundled/session-memory/handler.ts
  Homepage: https://docs.openclaw.ai/automation/hooks#session-memory
  Events: command:new

Requirements:
  Config: ✓ workspace.dir
\end{lstlisting}


\subsection{检查钩子资格}


\begin{lstlisting}[language=bash]
openclaw hooks check
\end{lstlisting}


显示钩子资格状态摘要（有多少已就绪，有多少未就绪）。

\textbf{选项：}

\begin{itemize}
  \item \texttt{--json}：以 JSON 格式输出
\end{itemize}


\textbf{示例输出：}

\begin{lstlisting}
Hooks Status

Total hooks: 4
Ready: 4
Not ready: 0
\end{lstlisting}


\subsection{启用钩子}


\begin{lstlisting}[language=bash]
openclaw hooks enable <name>
\end{lstlisting}


通过将特定钩子添加到配置（\texttt{\textasciitilde{}/.openclaw/config.json}）来启用它。

\textbf{注意：} 由插件管理的钩子在 \texttt{openclaw hooks list} 中显示 \texttt{plugin:}，
无法在此处启用/禁用。请改为启用/禁用该插件。

\textbf{参数：}

\begin{itemize}
  \item `\texttt{：钩子名称（例如 }session-memory`）
\end{itemize}


\textbf{示例：}

\begin{lstlisting}[language=bash]
openclaw hooks enable session-memory
\end{lstlisting}


\textbf{输出：}

\begin{lstlisting}
✓ Enabled hook: 💾 session-memory
\end{lstlisting}


\textbf{执行操作：}

\begin{itemize}
  \item 检查钩子是否存在且符合条件
  \item 在配置中更新 \texttt{hooks.internal.entries..enabled = true}
  \item 将配置保存到磁盘
\end{itemize}


\textbf{启用后：}

\begin{itemize}
  \item 重启 Gateway 网关以重新加载钩子（macOS 上重启菜单栏应用，或在开发环境中重启 Gateway 网关进程）。
\end{itemize}


\subsection{禁用钩子}


\begin{lstlisting}[language=bash]
openclaw hooks disable <name>
\end{lstlisting}


通过更新配置来禁用特定钩子。

\textbf{参数：}

\begin{itemize}
  \item `\texttt{：钩子名称（例如 }command-logger`）
\end{itemize}


\textbf{示例：}

\begin{lstlisting}[language=bash]
openclaw hooks disable command-logger
\end{lstlisting}


\textbf{输出：}

\begin{lstlisting}
⏸ Disabled hook: 📝 command-logger
\end{lstlisting}


\textbf{禁用后：}

\begin{itemize}
  \item 重启 Gateway 网关以重新加载钩子
\end{itemize}


\subsection{安装钩子}


\begin{lstlisting}[language=bash]
openclaw hooks install <path-or-spec>
\end{lstlisting}


从本地文件夹/压缩包或 npm 安装钩子包。

\textbf{执行操作：}

\begin{itemize}
  \item 将钩子包复制到 \texttt{\textasciitilde{}/.openclaw/hooks/}
  \item 在 \texttt{hooks.internal.entries.*} 中启用已安装的钩子
  \item 在 \texttt{hooks.internal.installs} 下记录安装信息
\end{itemize}


\textbf{选项：}

\begin{itemize}
  \item \texttt{-l, --link}：链接本地目录而不是复制（将其添加到 \texttt{hooks.internal.load.extraDirs}）
\end{itemize}


\textbf{支持的压缩包格式：} \texttt{.zip}、\texttt{.tgz}、\texttt{.tar.gz}、\texttt{.tar}

\textbf{示例：}

\begin{lstlisting}[language=bash]
# 本地目录
openclaw hooks install ./my-hook-pack

# 本地压缩包
openclaw hooks install ./my-hook-pack.zip

# NPM 包
openclaw hooks install @openclaw/my-hook-pack

# 链接本地目录而不复制
openclaw hooks install -l ./my-hook-pack
\end{lstlisting}


\subsection{更新钩子}


\begin{lstlisting}[language=bash]
openclaw hooks update <id>
openclaw hooks update --all
\end{lstlisting}


更新已安装的钩子包（仅限 npm 安装）。

\textbf{选项：}

\begin{itemize}
  \item \texttt{--all}：更新所有已跟踪的钩子包
  \item \texttt{--dry-run}：显示将要进行的更改，但不写入
\end{itemize}


\subsection{内置钩子}


\subsubsection{session-memory}


在你执行 \texttt{/new} 时将会话上下文保存到记忆中。

\textbf{启用：}

\begin{lstlisting}[language=bash]
openclaw hooks enable session-memory
\end{lstlisting}


\textbf{输出：} \texttt{\textasciitilde{}/.openclaw/workspace/memory/YYYY-MM-DD-slug.md}

\textbf{参见：} \href{/automation/hooks\#session-memory}{session-memory 文档}

\subsubsection{command-logger}


将所有命令事件记录到集中的审计文件中。

\textbf{启用：}

\begin{lstlisting}[language=bash]
openclaw hooks enable command-logger
\end{lstlisting}


\textbf{输出：} \texttt{\textasciitilde{}/.openclaw/logs/commands.log}

\textbf{查看日志：}

\begin{lstlisting}[language=bash]
# 最近的命令
tail -n 20 ~/.openclaw/logs/commands.log

# 格式化输出
cat ~/.openclaw/logs/commands.log | jq .

# 按操作过滤
grep '"action":"new"' ~/.openclaw/logs/commands.log | jq .
\end{lstlisting}


\textbf{参见：} \href{/automation/hooks\#command-logger}{command-logger 文档}

\subsubsection{boot-md}


在 Gateway 网关启动时（渠道启动后）运行 \texttt{BOOT.md}。

\textbf{事件}：\texttt{gateway:startup}

\textbf{启用}：

\begin{lstlisting}[language=bash]
openclaw hooks enable boot-md
\end{lstlisting}


\textbf{参见：} \href{/automation/hooks\#boot-md}{boot-md 文档}


\clearpage

% === 源文件: cli/index.md ===
\section{CLI 参考}


本页描述当前的 CLI 行为。如果命令发生变化，请更新此文档。

\subsection{命令页面}


\begin{itemize}
  \item \href{/cli/setup}{\texttt{setup}}
  \item \href{/cli/onboard}{\texttt{onboard}}
  \item \href{/cli/configure}{\texttt{configure}}
  \item \href{/cli/config}{\texttt{config}}
  \item \href{/cli/doctor}{\texttt{doctor}}
  \item \href{/cli/dashboard}{\texttt{dashboard}}
  \item \href{/cli/reset}{\texttt{reset}}
  \item \href{/cli/uninstall}{\texttt{uninstall}}
  \item \href{/cli/update}{\texttt{update}}
  \item \href{/cli/message}{\texttt{message}}
  \item \href{/cli/agent}{\texttt{agent}}
  \item \href{/cli/agents}{\texttt{agents}}
  \item \href{/cli/acp}{\texttt{acp}}
  \item \href{/cli/status}{\texttt{status}}
  \item \href{/cli/health}{\texttt{health}}
  \item \href{/cli/sessions}{\texttt{sessions}}
  \item \href{/cli/gateway}{\texttt{gateway}}
  \item \href{/cli/logs}{\texttt{logs}}
  \item \href{/cli/system}{\texttt{system}}
  \item \href{/cli/models}{\texttt{models}}
  \item \href{/cli/memory}{\texttt{memory}}
  \item \href{/cli/nodes}{\texttt{nodes}}
  \item \href{/cli/devices}{\texttt{devices}}
  \item \href{/cli/node}{\texttt{node}}
  \item \href{/cli/approvals}{\texttt{approvals}}
  \item \href{/cli/sandbox}{\texttt{sandbox}}
  \item \href{/cli/tui}{\texttt{tui}}
  \item \href{/cli/browser}{\texttt{browser}}
  \item \href{/cli/cron}{\texttt{cron}}
  \item \href{/cli/dns}{\texttt{dns}}
  \item \href{/cli/docs}{\texttt{docs}}
  \item \href{/cli/hooks}{\texttt{hooks}}
  \item \href{/cli/webhooks}{\texttt{webhooks}}
  \item \href{/cli/pairing}{\texttt{pairing}}
  \item \href{/cli/plugins}{\texttt{plugins}}（插件命令）
  \item \href{/cli/channels}{\texttt{channels}}
  \item \href{/cli/security}{\texttt{security}}
  \item \href{/cli/skills}{\texttt{skills}}
  \item \href{/cli/voicecall}{\texttt{voicecall}}（插件；如已安装）
\end{itemize}


\subsection{全局标志}


\begin{itemize}
  \item \texttt{--dev}：将状态隔离到 \texttt{\textasciitilde{}/.openclaw-dev} 下并调整默认端口。
  \item \texttt{--profile }：将状态隔离到 \texttt{\textasciitilde{}/.openclaw-} 下。
  \item \texttt{--no-color}：禁用 ANSI 颜色。
  \item \texttt{--update}：\texttt{openclaw update} 的简写（仅限源码安装）。
  \item \texttt{-V}、\texttt{--version}、\texttt{-v}：打印版本并退出。
\end{itemize}


\subsection{输出样式}


\begin{itemize}
  \item ANSI 颜色和进度指示器仅在 TTY 会话中渲染。
  \item OSC-8 超链接在支持的终端中渲染为可点击链接；否则回退到纯 URL。
  \item \texttt{--json}（以及支持的地方使用 \texttt{--plain}）禁用样式以获得干净输出。
  \item \texttt{--no-color} 禁用 ANSI 样式；也支持 \texttt{NO\_COLOR=1}。
  \item 长时间运行的命令显示进度指示器（支持时使用 OSC 9;4）。
\end{itemize}


\subsection{颜色调色板}


OpenClaw 在 CLI 输出中使用龙虾调色板。

\begin{itemize}
  \item \texttt{accent}（\#FF5A2D）：标题、标签、主要高亮。
  \item \texttt{accentBright}（\#FF7A3D）：命令名称、强调。
  \item \texttt{accentDim}（\#D14A22）：次要高亮文本。
  \item \texttt{info}（\#FF8A5B）：信息性值。
  \item \texttt{success}（\#2FBF71）：成功状态。
  \item \texttt{warn}（\#FFB020）：警告、回退、注意。
  \item \texttt{error}（\#E23D2D）：错误、失败。
  \item \texttt{muted}（\#8B7F77）：弱化、元数据。
\end{itemize}


调色板权威来源：\texttt{src/terminal/palette.ts}（又名"lobster seam"）。

\subsection{命令树}


\begin{lstlisting}
openclaw [--dev] [--profile <name>] <command>
  setup
  onboard
  configure
  config
    get
    set
    unset
  doctor
  security
    audit
  reset
  uninstall
  update
  channels
    list
    status
    logs
    add
    remove
    login
    logout
  skills
    list
    info
    check
  plugins
    list
    info
    install
    enable
    disable
    doctor
  memory
    status
    index
    search
  message
  agent
  agents
    list
    add
    delete
  acp
  status
  health
  sessions
  gateway
    call
    health
    status
    probe
    discover
    install
    uninstall
    start
    stop
    restart
    run
  logs
  system
    event
    heartbeat last|enable|disable
    presence
  models
    list
    status
    set
    set-image
    aliases list|add|remove
    fallbacks list|add|remove|clear
    image-fallbacks list|add|remove|clear
    scan
    auth add|setup-token|paste-token
    auth order get|set|clear
  sandbox
    list
    recreate
    explain
  cron
    status
    list
    add
    edit
    rm
    enable
    disable
    runs
    run
  nodes
  devices
  node
    run
    status
    install
    uninstall
    start
    stop
    restart
  approvals
    get
    set
    allowlist add|remove
  browser
    status
    start
    stop
    reset-profile
    tabs
    open
    focus
    close
    profiles
    create-profile
    delete-profile
    screenshot
    snapshot
    navigate
    resize
    click
    type
    press
    hover
    drag
    select
    upload
    fill
    dialog
    wait
    evaluate
    console
    pdf
  hooks
    list
    info
    check
    enable
    disable
    install
    update
  webhooks
    gmail setup|run
  pairing
    list
    approve
  docs
  dns
    setup
  tui
\end{lstlisting}


注意：插件可以添加额外的顶级命令（例如 \texttt{openclaw voicecall}）。

\subsection{安全}


\begin{itemize}
  \item \texttt{openclaw security audit} — 审计配置 + 本地状态中常见的安全隐患。
  \item \texttt{openclaw security audit --deep} — 尽力进行实时 Gateway 网关探测。
  \item \texttt{openclaw security audit --fix} — 收紧安全默认值并 chmod 状态/配置。
\end{itemize}


\subsection{插件}


管理扩展及其配置：

\begin{itemize}
  \item \texttt{openclaw plugins list} — 发现插件（使用 \texttt{--json} 获取机器可读输出）。
  \item \texttt{openclaw plugins info } — 显示插件详情。
  \item \texttt{openclaw plugins install } — 安装插件（或将插件路径添加到 \texttt{plugins.load.paths}）。
  \item \texttt{openclaw plugins enable } / \texttt{disable } — 切换 \texttt{plugins.entries..enabled}。
  \item \texttt{openclaw plugins doctor} — 报告插件加载错误。
\end{itemize}


大多数插件更改需要重启 Gateway 网关。参见 \href{/tools/plugin}{/plugin}。

\subsection{记忆}


对 \texttt{MEMORY.md} + \texttt{memory/*.md} 进行向量搜索：

\begin{itemize}
  \item \texttt{openclaw memory status} — 显示索引统计。
  \item \texttt{openclaw memory index} — 重新索引记忆文件。
  \item \texttt{openclaw memory search ""} — 对记忆进行语义搜索。
\end{itemize}


\subsection{聊天斜杠命令}


聊天消息支持 \texttt{/...} 命令（文本和原生）。参见 \href{/tools/slash-commands}{/tools/slash-commands}。

亮点：

\begin{itemize}
  \item \texttt{/status} 用于快速诊断。
  \item \texttt{/config} 用于持久化配置更改。
  \item \texttt{/debug} 用于仅运行时的配置覆盖（内存中，不写入磁盘；需要 \texttt{commands.debug: true}）。
\end{itemize}


\subsection{设置 + 新手引导}


\subsubsection{setup}


初始化配置 + 工作区。

选项：

\begin{itemize}
  \item \texttt{--workspace }：智能体工作区路径（默认 \texttt{\textasciitilde{}/.openclaw/workspace}）。
  \item \texttt{--wizard}：运行新手引导向导。
  \item \texttt{--non-interactive}：无提示运行向导。
  \item \texttt{--mode }：向导模式。
  \item \texttt{--remote-url }：远程 Gateway 网关 URL。
  \item \texttt{--remote-token }：远程 Gateway 网关令牌。
\end{itemize}


当存在任何向导标志（\texttt{--non-interactive}、\texttt{--mode}、\texttt{--remote-url}、\texttt{--remote-token}）时，向导自动运行。

\subsubsection{onboard}


交互式向导，用于设置 Gateway 网关、工作区和 Skills。

选项：

\begin{itemize}
  \item \texttt{--workspace }
  \item \texttt{--reset}（在向导之前重置配置 + 凭证 + 会话 + 工作区）
  \item \texttt{--non-interactive}
  \item \texttt{--mode }
  \item \texttt{--flow }（manual 是 advanced 的别名）
  \item \texttt{--auth-choice }
  \item \texttt{--token-provider }（非交互式；与 \texttt{--auth-choice token} 配合使用）
  \item \texttt{--token }（非交互式；与 \texttt{--auth-choice token} 配合使用）
  \item \texttt{--token-profile-id }（非交互式；默认：\texttt{:manual}）
  \item \texttt{--token-expires-in }（非交互式；例如 \texttt{365d}、\texttt{12h}）
  \item \texttt{--anthropic-api-key }
  \item \texttt{--openai-api-key }
  \item \texttt{--openrouter-api-key }
  \item \texttt{--ai-gateway-api-key }
  \item \texttt{--moonshot-api-key }
  \item \texttt{--kimi-code-api-key }
  \item \texttt{--gemini-api-key }
  \item \texttt{--zai-api-key }
  \item \texttt{--minimax-api-key }
  \item \texttt{--opencode-zen-api-key }
  \item \texttt{--gateway-port }
  \item \texttt{--gateway-bind }
  \item \texttt{--gateway-auth }
  \item \texttt{--gateway-token }
  \item \texttt{--gateway-password }
  \item \texttt{--remote-url }
  \item \texttt{--remote-token }
  \item \texttt{--tailscale }
  \item \texttt{--tailscale-reset-on-exit}
  \item \texttt{--install-daemon}
  \item \texttt{--no-install-daemon}（别名：\texttt{--skip-daemon}）
  \item \texttt{--daemon-runtime }
  \item \texttt{--skip-channels}
  \item \texttt{--skip-skills}
  \item \texttt{--skip-health}
  \item \texttt{--skip-ui}
  \item \texttt{--node-manager }（推荐 pnpm；不建议将 bun 用于 Gateway 网关运行时）
  \item \texttt{--json}
\end{itemize}


\subsubsection{configure}


交互式配置向导（模型、渠道、Skills、Gateway 网关）。

\subsubsection{config}


非交互式配置辅助工具（get/set/unset）。不带子命令运行 \texttt{openclaw config} 会启动向导。

子命令：

\begin{itemize}
  \item \texttt{config get }：打印配置值（点/括号路径）。
  \item \texttt{config set  }：设置值（JSON5 或原始字符串）。
  \item \texttt{config unset }：删除值。
\end{itemize}


\subsubsection{doctor}


健康检查 + 快速修复（配置 + Gateway 网关 + 旧版服务）。

选项：

\begin{itemize}
  \item \texttt{--no-workspace-suggestions}：禁用工作区记忆提示。
  \item \texttt{--yes}：无提示接受默认值（无头模式）。
  \item \texttt{--non-interactive}：跳过提示；仅应用安全迁移。
  \item \texttt{--deep}：扫描系统服务以查找额外的 Gateway 网关安装。
\end{itemize}


\subsection{渠道辅助工具}


\subsubsection{channels}


管理聊天渠道账户（WhatsApp/Telegram/Discord/Google Chat/Slack/Mattermost（插件）/Signal/iMessage/MS Teams）。

子命令：

\begin{itemize}
  \item \texttt{channels list}：显示已配置的渠道和认证配置文件。
  \item \texttt{channels status}：检查 Gateway 网关可达性和渠道健康状况（\texttt{--probe} 运行额外检查；使用 \texttt{openclaw health} 或 \texttt{openclaw status --deep} 进行 Gateway 网关健康探测）。
  \item 提示：\texttt{channels status} 在检测到常见配置错误时会打印带有建议修复的警告（然后指向 \texttt{openclaw doctor}）。
  \item \texttt{channels logs}：显示 Gateway 网关日志文件中最近的渠道日志。
  \item \texttt{channels add}：不传标志时使用向导式设置；标志切换到非交互模式。
  \item \texttt{channels remove}：默认禁用；传 \texttt{--delete} 可无提示删除配置条目。
  \item \texttt{channels login}：交互式渠道登录（仅限 WhatsApp Web）。
  \item \texttt{channels logout}：登出渠道会话（如支持）。
\end{itemize}


通用选项：

\begin{itemize}
  \item \texttt{--channel }：\texttt{whatsapp|telegram|discord|googlechat|slack|mattermost|signal|imessage|msteams}
  \item \texttt{--account }：渠道账户 id（默认 \texttt{default}）
  \item \texttt{--name }：账户的显示名称
\end{itemize}


\texttt{channels login} 选项：

\begin{itemize}
  \item \texttt{--channel }（默认 \texttt{whatsapp}；支持 \texttt{whatsapp}/\texttt{web}）
  \item \texttt{--account }
  \item \texttt{--verbose}
\end{itemize}


\texttt{channels logout} 选项：

\begin{itemize}
  \item \texttt{--channel }（默认 \texttt{whatsapp}）
  \item \texttt{--account }
\end{itemize}


\texttt{channels list} 选项：

\begin{itemize}
  \item \texttt{--no-usage}：跳过模型提供商用量/配额快照（仅限 OAuth/API 支持的）。
  \item \texttt{--json}：输出 JSON（除非设置 \texttt{--no-usage}，否则包含用量）。
\end{itemize}


\texttt{channels logs} 选项：

\begin{itemize}
  \item \texttt{--channel }（默认 \texttt{all}）
  \item \texttt{--lines }（默认 \texttt{200}）
  \item \texttt{--json}
\end{itemize}


更多详情：\href{/concepts/oauth}{/concepts/oauth}

示例：

\begin{lstlisting}[language=bash]
openclaw channels add --channel telegram --account alerts --name "Alerts Bot" --token $TELEGRAM_BOT_TOKEN
openclaw channels add --channel discord --account work --name "Work Bot" --token $DISCORD_BOT_TOKEN
openclaw channels remove --channel discord --account work --delete
openclaw channels status --probe
openclaw status --deep
\end{lstlisting}


\subsubsection{skills}


列出和检查可用的 Skills 及就绪信息。

子命令：

\begin{itemize}
  \item \texttt{skills list}：列出 Skills（无子命令时的默认行为）。
  \item \texttt{skills info }：显示单个 Skill 的详情。
  \item \texttt{skills check}：就绪与缺失需求的摘要。
\end{itemize}


选项：

\begin{itemize}
  \item \texttt{--eligible}：仅显示就绪的 Skills。
  \item \texttt{--json}：输出 JSON（无样式）。
  \item \texttt{-v}、\texttt{--verbose}：包含缺失需求详情。
\end{itemize}


提示：使用 \texttt{npx clawhub} 搜索、安装和同步 Skills。

\subsubsection{pairing}


批准跨渠道的私信配对请求。

子命令：

\begin{itemize}
  \item \texttt{pairing list  [--json]}
  \item \texttt{pairing approve   [--notify]}
\end{itemize}


\subsubsection{webhooks gmail}


Gmail Pub/Sub 钩子设置 + 运行器。参见 \href{/automation/gmail-pubsub}{/automation/gmail-pubsub}。

子命令：

\begin{itemize}
  \item \texttt{webhooks gmail setup}（需要 \texttt{--account }；支持 \texttt{--project}、\texttt{--topic}、\texttt{--subscription}、\texttt{--label}、\texttt{--hook-url}、\texttt{--hook-token}、\texttt{--push-token}、\texttt{--bind}、\texttt{--port}、\texttt{--path}、\texttt{--include-body}、\texttt{--max-bytes}、\texttt{--renew-minutes}、\texttt{--tailscale}、\texttt{--tailscale-path}、\texttt{--tailscale-target}、\texttt{--push-endpoint}、\texttt{--json}）
  \item \texttt{webhooks gmail run}（相同标志的运行时覆盖）
\end{itemize}


\subsubsection{dns setup}


广域发现 DNS 辅助工具（CoreDNS + Tailscale）。参见 \href{/gateway/discovery}{/gateway/discovery}。

选项：

\begin{itemize}
  \item \texttt{--apply}：安装/更新 CoreDNS 配置（需要 sudo；仅限 macOS）。
\end{itemize}


\subsection{消息 + 智能体}


\subsubsection{message}


统一的出站消息 + 渠道操作。

参见：\href{/cli/message}{/cli/message}

子命令：

\begin{itemize}
  \item \texttt{message send|poll|react|reactions|read|edit|delete|pin|unpin|pins|permissions|search|timeout|kick|ban}
  \item \texttt{message thread }
  \item \texttt{message emoji }
  \item \texttt{message sticker }
  \item \texttt{message role }
  \item \texttt{message channel }
  \item \texttt{message member info}
  \item \texttt{message voice status}
  \item \texttt{message event }
\end{itemize}


示例：

\begin{itemize}
  \item \texttt{openclaw message send --target +15555550123 --message "Hi"}
  \item \texttt{openclaw message poll --channel discord --target channel:123 --poll-question "Snack?" --poll-option Pizza --poll-option Sushi}
\end{itemize}


\subsubsection{agent}


通过 Gateway 网关运行一个智能体回合（或使用 \texttt{--local} 嵌入式运行）。

必需：

\begin{itemize}
  \item \texttt{--message }
\end{itemize}


选项：

\begin{itemize}
  \item \texttt{--to }（用于会话键和可选发送）
  \item \texttt{--session-id }
  \item \texttt{--thinking }（仅限 GPT-5.2 + Codex 模型）
  \item \texttt{--verbose }
  \item \texttt{--channel }
  \item \texttt{--local}
  \item \texttt{--deliver}
  \item \texttt{--json}
  \item \texttt{--timeout }
\end{itemize}


\subsubsection{agents}


管理隔离的智能体（工作区 + 认证 + 路由）。

\paragraph{agents list}


列出已配置的智能体。

选项：

\begin{itemize}
  \item \texttt{--json}
  \item \texttt{--bindings}
\end{itemize}


\paragraph{agents add [name]}


添加新的隔离智能体。除非传入标志（或 \texttt{--non-interactive}），否则运行引导向导；非交互模式下 \texttt{--workspace} 是必需的。

选项：

\begin{itemize}
  \item \texttt{--workspace }
  \item \texttt{--model }
  \item \texttt{--agent-dir }
  \item \texttt{--bind }（可重复）
  \item \texttt{--non-interactive}
  \item \texttt{--json}
\end{itemize}


绑定规范使用 \texttt{channel[:accountId]}。对于 WhatsApp，省略 \texttt{accountId} 时使用默认账户 id。

\paragraph{agents delete }


删除智能体并清理其工作区 + 状态。

选项：

\begin{itemize}
  \item \texttt{--force}
  \item \texttt{--json}
\end{itemize}


\subsubsection{acp}


运行连接 IDE 到 Gateway 网关的 ACP 桥接。

完整选项和示例参见 \href{/cli/acp}{\texttt{acp}}。

\subsubsection{status}


显示关联会话健康状况和最近的收件人。

选项：

\begin{itemize}
  \item \texttt{--json}
  \item \texttt{--all}（完整诊断；只读，可粘贴）
  \item \texttt{--deep}（探测渠道）
  \item \texttt{--usage}（显示模型提供商用量/配额）
  \item \texttt{--timeout }
  \item \texttt{--verbose}
  \item \texttt{--debug}（\texttt{--verbose} 的别名）
\end{itemize}


说明：

\begin{itemize}
  \item 概览包含 Gateway 网关 + 节点主机服务状态（如可用）。
\end{itemize}


\subsubsection{用量跟踪}


当 OAuth/API 凭证可用时，OpenClaw 可以显示提供商用量/配额。

显示位置：

\begin{itemize}
  \item \texttt{/status}（可用时添加简短的提供商用量行）
  \item \texttt{openclaw status --usage}（打印完整的提供商明细）
  \item macOS 菜单栏（上下文下的用量部分）
\end{itemize}


说明：

\begin{itemize}
  \item 数据直接来自提供商用量端点（非估算）。
  \item 提供商：Anthropic、GitHub Copilot、OpenAI Codex OAuth，以及启用这些提供商插件时的 Gemini CLI/Antigravity。
  \item 如果没有匹配的凭证，用量会被隐藏。
  \item 详情：参见\href{/concepts/usage-tracking}{用量跟踪}。
\end{itemize}


\subsubsection{health}


从运行中的 Gateway 网关获取健康状态。

选项：

\begin{itemize}
  \item \texttt{--json}
  \item \texttt{--timeout }
  \item \texttt{--verbose}
\end{itemize}


\subsubsection{sessions}


列出存储的对话会话。

选项：

\begin{itemize}
  \item \texttt{--json}
  \item \texttt{--verbose}
  \item \texttt{--store }
  \item \texttt{--active }
\end{itemize}


\subsection{重置/卸载}


\subsubsection{reset}


重置本地配置/状态（保留 CLI 安装）。

选项：

\begin{itemize}
  \item \texttt{--scope }
  \item \texttt{--yes}
  \item \texttt{--non-interactive}
  \item \texttt{--dry-run}
\end{itemize}


说明：

\begin{itemize}
  \item \texttt{--non-interactive} 需要 \texttt{--scope} 和 \texttt{--yes}。
\end{itemize}


\subsubsection{uninstall}


卸载 Gateway 网关服务 + 本地数据（CLI 保留）。

选项：

\begin{itemize}
  \item \texttt{--service}
  \item \texttt{--state}
  \item \texttt{--workspace}
  \item \texttt{--app}
  \item \texttt{--all}
  \item \texttt{--yes}
  \item \texttt{--non-interactive}
  \item \texttt{--dry-run}
\end{itemize}


说明：

\begin{itemize}
  \item \texttt{--non-interactive} 需要 \texttt{--yes} 和明确的范围（或 \texttt{--all}）。
\end{itemize}


\subsection{Gateway 网关}


\subsubsection{gateway}


运行 WebSocket Gateway 网关。

选项：

\begin{itemize}
  \item \texttt{--port }
  \item \texttt{--bind }
  \item \texttt{--token }
  \item \texttt{--auth }
  \item \texttt{--password }
  \item \texttt{--tailscale }
  \item \texttt{--tailscale-reset-on-exit}
  \item \texttt{--allow-unconfigured}
  \item \texttt{--dev}
  \item \texttt{--reset}（重置 dev 配置 + 凭证 + 会话 + 工作区）
  \item \texttt{--force}（终止端口上的现有监听器）
  \item \texttt{--verbose}
  \item \texttt{--claude-cli-logs}
  \item \texttt{--ws-log }
  \item \texttt{--compact}（\texttt{--ws-log compact} 的别名）
  \item \texttt{--raw-stream}
  \item \texttt{--raw-stream-path }
\end{itemize}


\subsubsection{gateway service}


管理 Gateway 网关服务（launchd/systemd/schtasks）。

子命令：

\begin{itemize}
  \item \texttt{gateway status}（默认探测 Gateway 网关 RPC）
  \item \texttt{gateway install}（服务安装）
  \item \texttt{gateway uninstall}
  \item \texttt{gateway start}
  \item \texttt{gateway stop}
  \item \texttt{gateway restart}
\end{itemize}


说明：

\begin{itemize}
  \item \texttt{gateway status} 默认使用服务解析的端口/配置探测 Gateway 网关 RPC（使用 \texttt{--url/--token/--password} 覆盖）。
  \item \texttt{gateway status} 支持 \texttt{--no-probe}、\texttt{--deep} 和 \texttt{--json} 用于脚本化。
  \item \texttt{gateway status} 在检测到旧版或额外的 Gateway 网关服务时也会显示（\texttt{--deep} 添加系统级扫描）。配置文件命名的 OpenClaw 服务被视为一等公民，不会被标记为"额外"。
  \item \texttt{gateway status} 打印 CLI 使用的配置路径与服务可能使用的配置（服务环境），以及解析的探测目标 URL。
  \item \texttt{gateway install|uninstall|start|stop|restart} 支持 \texttt{--json} 用于脚本化（默认输出保持人类友好）。
  \item \texttt{gateway install} 默认使用 Node 运行时；\textbf{不建议}使用 bun（WhatsApp/Telegram bug）。
  \item \texttt{gateway install} 选项：\texttt{--port}、\texttt{--runtime}、\texttt{--token}、\texttt{--force}、\texttt{--json}。
\end{itemize}


\subsubsection{logs}


通过 RPC 跟踪 Gateway 网关文件日志。

说明：

\begin{itemize}
  \item TTY 会话渲染彩色、结构化视图；非 TTY 回退到纯文本。
  \item \texttt{--json} 输出行分隔的 JSON（每行一个日志事件）。
\end{itemize}


示例：

\begin{lstlisting}[language=bash]
openclaw logs --follow
openclaw logs --limit 200
openclaw logs --plain
openclaw logs --json
openclaw logs --no-color
\end{lstlisting}


\subsubsection{gateway }


Gateway 网关 CLI 辅助工具（RPC 子命令使用 \texttt{--url}、\texttt{--token}、\texttt{--password}、\texttt{--timeout}、\texttt{--expect-final}）。

子命令：

\begin{itemize}
  \item \texttt{gateway call  [--params ]}
  \item \texttt{gateway health}
  \item \texttt{gateway status}
  \item \texttt{gateway probe}
  \item \texttt{gateway discover}
  \item \texttt{gateway install|uninstall|start|stop|restart}
  \item \texttt{gateway run}
\end{itemize}


常见 RPC：

\begin{itemize}
  \item \texttt{config.apply}（验证 + 写入配置 + 重启 + 唤醒）
  \item \texttt{config.patch}（合并部分更新 + 重启 + 唤醒）
  \item \texttt{update.run}（运行更新 + 重启 + 唤醒）
\end{itemize}


提示：直接调用 \texttt{config.set}/\texttt{config.apply}/\texttt{config.patch} 时，如果配置已存在，请传入来自 \texttt{config.get} 的 \texttt{baseHash}。

\subsection{模型}


回退行为和扫描策略参见 \href{/concepts/models}{/concepts/models}。

首选 Anthropic 认证（setup-token）：

\begin{lstlisting}[language=bash]
claude setup-token
openclaw models auth setup-token --provider anthropic
openclaw models status
\end{lstlisting}


\subsubsection{models（根命令）}


\texttt{openclaw models} 是 \texttt{models status} 的别名。

根选项：

\begin{itemize}
  \item \texttt{--status-json}（\texttt{models status --json} 的别名）
  \item \texttt{--status-plain}（\texttt{models status --plain} 的别名）
\end{itemize}


\subsubsection{models list}


选项：

\begin{itemize}
  \item \texttt{--all}
  \item \texttt{--local}
  \item \texttt{--provider }
  \item \texttt{--json}
  \item \texttt{--plain}
\end{itemize}


\subsubsection{models status}


选项：

\begin{itemize}
  \item \texttt{--json}
  \item \texttt{--plain}
  \item \texttt{--check}（退出码 1=过期/缺失，2=即将过期）
  \item \texttt{--probe}（对已配置认证配置文件进行实时探测）
  \item \texttt{--probe-provider }
  \item \texttt{--probe-profile }（重复或逗号分隔）
  \item \texttt{--probe-timeout }
  \item \texttt{--probe-concurrency }
  \item \texttt{--probe-max-tokens }
\end{itemize}


始终包含认证概览和认证存储中配置文件的 OAuth 过期状态。\texttt{--probe} 运行实时请求（可能消耗令牌并触发速率限制）。

\subsubsection{models set }


设置 \texttt{agents.defaults.model.primary}。

\subsubsection{models set-image }


设置 \texttt{agents.defaults.imageModel.primary}。

\subsubsection{models aliases list|add|remove}


选项：

\begin{itemize}
  \item \texttt{list}：\texttt{--json}、\texttt{--plain}
  \item \texttt{add  }
  \item \texttt{remove }
\end{itemize}


\subsubsection{models fallbacks list|add|remove|clear}


选项：

\begin{itemize}
  \item \texttt{list}：\texttt{--json}、\texttt{--plain}
  \item \texttt{add }
  \item \texttt{remove }
  \item \texttt{clear}
\end{itemize}


\subsubsection{models image-fallbacks list|add|remove|clear}


选项：

\begin{itemize}
  \item \texttt{list}：\texttt{--json}、\texttt{--plain}
  \item \texttt{add }
  \item \texttt{remove }
  \item \texttt{clear}
\end{itemize}


\subsubsection{models scan}


选项：

\begin{itemize}
  \item \texttt{--min-params }
  \item \texttt{--max-age-days }
  \item \texttt{--provider }
  \item \texttt{--max-candidates }
  \item \texttt{--timeout }
  \item \texttt{--concurrency }
  \item \texttt{--no-probe}
  \item \texttt{--yes}
  \item \texttt{--no-input}
  \item \texttt{--set-default}
  \item \texttt{--set-image}
  \item \texttt{--json}
\end{itemize}


\subsubsection{models auth add|setup-token|paste-token}


选项：

\begin{itemize}
  \item \texttt{add}：交互式认证辅助工具
  \item \texttt{setup-token}：\texttt{--provider }（默认 \texttt{anthropic}）、\texttt{--yes}
  \item \texttt{paste-token}：\texttt{--provider }、\texttt{--profile-id }、\texttt{--expires-in }
\end{itemize}


\subsubsection{models auth order get|set|clear}


选项：

\begin{itemize}
  \item \texttt{get}：\texttt{--provider }、\texttt{--agent }、\texttt{--json}
  \item \texttt{set}：\texttt{--provider }、\texttt{--agent }、``
  \item \texttt{clear}：\texttt{--provider }、\texttt{--agent }
\end{itemize}


\subsection{系统}


\subsubsection{system event}


将系统事件加入队列并可选触发心跳（Gateway 网关 RPC）。

必需：

\begin{itemize}
  \item \texttt{--text }
\end{itemize}


选项：

\begin{itemize}
  \item \texttt{--mode }
  \item \texttt{--json}
  \item \texttt{--url}、\texttt{--token}、\texttt{--timeout}、\texttt{--expect-final}
\end{itemize}


\subsubsection{system heartbeat last|enable|disable}


心跳控制（Gateway 网关 RPC）。

选项：

\begin{itemize}
  \item \texttt{--json}
  \item \texttt{--url}、\texttt{--token}、\texttt{--timeout}、\texttt{--expect-final}
\end{itemize}


\subsubsection{system presence}


列出系统存在条目（Gateway 网关 RPC）。

选项：

\begin{itemize}
  \item \texttt{--json}
  \item \texttt{--url}、\texttt{--token}、\texttt{--timeout}、\texttt{--expect-final}
\end{itemize}


\subsection{定时任务}


管理计划任务（Gateway 网关 RPC）。参见 \href{/automation/cron-jobs}{/automation/cron-jobs}。

子命令：

\begin{itemize}
  \item \texttt{cron status [--json]}
  \item \texttt{cron list [--all] [--json]}（默认表格输出；使用 \texttt{--json} 获取原始数据）
  \item \texttt{cron add}（别名：\texttt{create}；需要 \texttt{--name} 和 \texttt{--at} | \texttt{--every} | \texttt{--cron} 三选一，以及 \texttt{--system-event} | \texttt{--message} 负载二选一）
  \item \texttt{cron edit }（补丁字段）
  \item \texttt{cron rm }（别名：\texttt{remove}、\texttt{delete}）
  \item \texttt{cron enable }
  \item \texttt{cron disable }
  \item \texttt{cron runs --id  [--limit ]}
  \item \texttt{cron run  [--force]}
\end{itemize}


所有 \texttt{cron} 命令接受 \texttt{--url}、\texttt{--token}、\texttt{--timeout}、\texttt{--expect-final}。

\subsection{节点主机}


\texttt{node} 运行\textbf{无头节点主机}或将其作为后台服务管理。参见 \href{/cli/node}{\texttt{openclaw node}}。

子命令：

\begin{itemize}
  \item \texttt{node run --host  --port 18789}
  \item \texttt{node status}
  \item \texttt{node install [--host ] [--port ] [--tls] [--tls-fingerprint ] [--node-id ] [--display-name ] [--runtime ] [--force]}
  \item \texttt{node uninstall}
  \item \texttt{node stop}
  \item \texttt{node restart}
\end{itemize}


\subsection{节点}


\texttt{nodes} 与 Gateway 网关通信并针对已配对的节点。参见 \href{/nodes}{/nodes}。

通用选项：

\begin{itemize}
  \item \texttt{--url}、\texttt{--token}、\texttt{--timeout}、\texttt{--json}
\end{itemize}


子命令：

\begin{itemize}
  \item \texttt{nodes status [--connected] [--last-connected ]}
  \item \texttt{nodes describe --node }
  \item \texttt{nodes list [--connected] [--last-connected ]}
  \item \texttt{nodes pending}
  \item \texttt{nodes approve }
  \item \texttt{nodes reject }
  \item \texttt{nodes rename --node  --name }
  \item \texttt{nodes invoke --node  --command  [--params ] [--invoke-timeout ] [--idempotency-key ]}
  \item \texttt{nodes run --node  [--cwd ] [--env KEY=VAL] [--command-timeout ] [--needs-screen-recording] [--invoke-timeout ] }（mac 节点或无头节点主机）
  \item \texttt{nodes notify --node  [--title ] [--body ] [--sound ] [--priority ] [--delivery ] [--invoke-timeout ]}（仅限 mac）
\end{itemize}


相机：

\begin{itemize}
  \item \texttt{nodes camera list --node }
  \item \texttt{nodes camera snap --node  [--facing front|back|both] [--device-id ] [--max-width ] [--quality <0-1>] [--delay-ms ] [--invoke-timeout ]}
  \item \texttt{nodes camera clip --node  [--facing front|back] [--device-id ] [--duration ] [--no-audio] [--invoke-timeout ]}
\end{itemize}


画布 + 屏幕：

\begin{itemize}
  \item \texttt{nodes canvas snapshot --node  [--format png|jpg|jpeg] [--max-width ] [--quality <0-1>] [--invoke-timeout ]}
  \item \texttt{nodes canvas present --node  [--target ] [--x ] [--y ] [--width ] [--height ] [--invoke-timeout ]}
  \item \texttt{nodes canvas hide --node  [--invoke-timeout ]}
  \item \texttt{nodes canvas navigate  --node  [--invoke-timeout ]}
  \item \texttt{nodes canvas eval [] --node  [--js ] [--invoke-timeout ]}
  \item \texttt{nodes canvas a2ui push --node  (--jsonl  | --text ) [--invoke-timeout ]}
  \item \texttt{nodes canvas a2ui reset --node  [--invoke-timeout ]}
  \item \texttt{nodes screen record --node  [--screen ] [--duration ] [--fps ] [--no-audio] [--out ] [--invoke-timeout ]}
\end{itemize}


位置：

\begin{itemize}
  \item \texttt{nodes location get --node  [--max-age ] [--accuracy ] [--location-timeout ] [--invoke-timeout ]}
\end{itemize}


\subsection{浏览器}


浏览器控制 CLI（专用 Chrome/Brave/Edge/Chromium）。参见 \href{/cli/browser}{\texttt{openclaw browser}} 和\href{/tools/browser}{浏览器工具}。

通用选项：

\begin{itemize}
  \item \texttt{--url}、\texttt{--token}、\texttt{--timeout}、\texttt{--json}
  \item \texttt{--browser-profile }
\end{itemize}


管理：

\begin{itemize}
  \item \texttt{browser status}
  \item \texttt{browser start}
  \item \texttt{browser stop}
  \item \texttt{browser reset-profile}
  \item \texttt{browser tabs}
  \item \texttt{browser open }
  \item \texttt{browser focus }
  \item \texttt{browser close [targetId]}
  \item \texttt{browser profiles}
  \item \texttt{browser create-profile --name  [--color ] [--cdp-url ]}
  \item \texttt{browser delete-profile --name }
\end{itemize}


检查：

\begin{itemize}
  \item \texttt{browser screenshot [targetId] [--full-page] [--ref ] [--element ] [--type png|jpeg]}
  \item \texttt{browser snapshot [--format aria|ai] [--target-id ] [--limit ] [--interactive] [--compact] [--depth ] [--selector ] [--out ]}
\end{itemize}


操作：

\begin{itemize}
  \item \texttt{browser navigate  [--target-id ]}
  \item \texttt{browser resize   [--target-id ]}
  \item \texttt{browser click  [--double] [--button ] [--modifiers ] [--target-id ]}
  \item \texttt{browser type   [--submit] [--slowly] [--target-id ]}
  \item \texttt{browser press  [--target-id ]}
  \item \texttt{browser hover  [--target-id ]}
  \item \texttt{browser drag   [--target-id ]}
  \item \texttt{browser select   [--target-id ]}
  \item \texttt{browser upload  [--ref ] [--input-ref ] [--element ] [--target-id ] [--timeout-ms ]}
  \item \texttt{browser fill [--fields ] [--fields-file ] [--target-id ]}
  \item \texttt{browser dialog --accept|--dismiss [--prompt ] [--target-id ] [--timeout-ms ]}
  \item \texttt{browser wait [--time ] [--text ] [--text-gone ] [--target-id ]}
  \item \texttt{browser evaluate --fn  [--ref ] [--target-id ]}
  \item \texttt{browser console [--level ] [--target-id ]}
  \item \texttt{browser pdf [--target-id ]}
\end{itemize}


\subsection{文档搜索}


\subsubsection{docs [query...]}


搜索在线文档索引。

\subsection{TUI}


\subsubsection{tui}


打开连接到 Gateway 网关的终端 UI。

选项：

\begin{itemize}
  \item \texttt{--url }
  \item \texttt{--token }
  \item \texttt{--password }
  \item \texttt{--session }
  \item \texttt{--deliver}
  \item \texttt{--thinking }
  \item \texttt{--message }
  \item \texttt{--timeout-ms }（默认为 \texttt{agents.defaults.timeoutSeconds}）
  \item \texttt{--history-limit }
\end{itemize}



\clearpage

% === 源文件: cli/logs.md ===
\section{openclaw logs}


通过 RPC 跟踪 Gateway 网关文件日志（在远程模式下可用）。

相关内容：

\begin{itemize}
  \item 日志概述：\href{/logging}{日志}
\end{itemize}


\subsection{示例}


\begin{lstlisting}[language=bash]
openclaw logs
openclaw logs --follow
openclaw logs --json
openclaw logs --limit 500
\end{lstlisting}



\clearpage

% === 源文件: cli/memory.md ===
\section{openclaw memory}


管理语义记忆的索引和搜索。
由活跃的记忆插件提供（默认：\texttt{memory-core}；设置 \texttt{plugins.slots.memory = "none"} 可禁用）。

相关内容：

\begin{itemize}
  \item 记忆概念：\href{/concepts/memory}{记忆}
  \item 插件：\href{/tools/plugin}{插件}
\end{itemize}


\subsection{示例}


\begin{lstlisting}[language=bash]
openclaw memory status
openclaw memory status --deep
openclaw memory status --deep --index
openclaw memory status --deep --index --verbose
openclaw memory index
openclaw memory index --verbose
openclaw memory search "release checklist"
openclaw memory status --agent main
openclaw memory index --agent main --verbose
\end{lstlisting}


\subsection{选项}


通用选项：

\begin{itemize}
  \item \texttt{--agent }：限定到单个智能体（默认：所有已配置的智能体）。
  \item \texttt{--verbose}：在探测和索引期间输出详细日志。
\end{itemize}


说明：

\begin{itemize}
  \item \texttt{memory status --deep} 探测向量存储和嵌入模型的可用性。
  \item \texttt{memory status --deep --index} 在存储有未同步变更时运行重新索引。
  \item \texttt{memory index --verbose} 打印每个阶段的详细信息（提供商、模型、数据源、批处理活动）。
  \item \texttt{memory status} 包含通过 \texttt{memorySearch.extraPaths} 配置的所有额外路径。
\end{itemize}



\clearpage

% === 源文件: cli/message.md ===
\section{openclaw message}


用于发送消息和渠道操作的单一出站命令
（Discord/Google Chat/Slack/Mattermost（插件）/Telegram/WhatsApp/Signal/iMessage/MS Teams）。

\subsection{用法}


\begin{lstlisting}
openclaw message <subcommand> [flags]
\end{lstlisting}


渠道选择：

\begin{itemize}
  \item 如果配置了多个渠道，则必须指定 \texttt{--channel}。
  \item 如果只配置了一个渠道，则该渠道为默认值。
  \item 可选值：\texttt{whatsapp|telegram|discord|googlechat|slack|mattermost|signal|imessage|msteams}（Mattermost 需要插件）
\end{itemize}


目标格式（\texttt{--target}）：

\begin{itemize}
  \item WhatsApp：E.164 或群组 JID
  \item Telegram：聊天 ID 或 \texttt{@username}
  \item Discord：\texttt{channel:} 或 \texttt{user:}（或 \texttt{<@id>} 提及；纯数字 ID 被视为频道）
  \item Google Chat：\texttt{spaces/} 或 \texttt{users/}
  \item Slack：\texttt{channel:} 或 \texttt{user:}（接受纯频道 ID）
  \item Mattermost（插件）：\texttt{channel:}、\texttt{user:} 或 \texttt{@username}（纯 ID 被视为频道）
  \item Signal：\texttt{+E.164}、\texttt{group:}、\texttt{signal:+E.164}、\texttt{signal:group:} 或 \texttt{username:}/\texttt{u:}
  \item iMessage：句柄、\texttt{chat\_id:}、\texttt{chat\_guid:} 或 \texttt{chat\_identifier:}
  \item MS Teams：会话 ID（\texttt{19:...@thread.tacv2}）或 \texttt{conversation:} 或 \texttt{user:}
\end{itemize}


名称查找：

\begin{itemize}
  \item 对于支持的提供商（Discord/Slack 等），如 \texttt{Help} 或 \texttt{\#help} 之类的频道名称会通过目录缓存进行解析。
  \item 缓存未命中时，如果提供商支持，OpenClaw 将尝试实时目录查找。
\end{itemize}


\subsection{通用标志}


\begin{itemize}
  \item \texttt{--channel }
  \item \texttt{--account }
  \item \texttt{--target }（用于 send/poll/read 等的目标渠道或用户）
  \item \texttt{--targets }（可重复；仅限广播）
  \item \texttt{--json}
  \item \texttt{--dry-run}
  \item \texttt{--verbose}
\end{itemize}


\subsection{操作}


\subsubsection{核心}


\begin{itemize}
  \item \texttt{send}
  \item 渠道：WhatsApp/Telegram/Discord/Google Chat/Slack/Mattermost（插件）/Signal/iMessage/MS Teams
  \item 必需：\texttt{--target}，以及 \texttt{--message} 或 \texttt{--media}
  \item 可选：\texttt{--media}、\texttt{--reply-to}、\texttt{--thread-id}、\texttt{--gif-playback}
  \item 仅限 Telegram：\texttt{--buttons}（需要 \texttt{channels.telegram.capabilities.inlineButtons} 以启用）
  \item 仅限 Telegram：\texttt{--thread-id}（论坛主题 ID）
  \item 仅限 Slack：\texttt{--thread-id}（线程时间戳；\texttt{--reply-to} 使用相同字段）
  \item 仅限 WhatsApp：\texttt{--gif-playback}
\end{itemize}


\begin{itemize}
  \item \texttt{poll}
  \item 渠道：WhatsApp/Discord/MS Teams
  \item 必需：\texttt{--target}、\texttt{--poll-question}、\texttt{--poll-option}（可重复）
  \item 可选：\texttt{--poll-multi}
  \item 仅限 Discord：\texttt{--poll-duration-hours}、\texttt{--message}
\end{itemize}


\begin{itemize}
  \item \texttt{react}
  \item 渠道：Discord/Google Chat/Slack/Telegram/WhatsApp/Signal
  \item 必需：\texttt{--message-id}、\texttt{--target}
  \item 可选：\texttt{--emoji}、\texttt{--remove}、\texttt{--participant}、\texttt{--from-me}、\texttt{--target-author}、\texttt{--target-author-uuid}
  \item 注意：\texttt{--remove} 需要 \texttt{--emoji}（省略 \texttt{--emoji} 可清除自己的表情回应（如果支持）；参见 /tools/reactions）
  \item 仅限 WhatsApp：\texttt{--participant}、\texttt{--from-me}
  \item Signal 群组表情回应：需要 \texttt{--target-author} 或 \texttt{--target-author-uuid}
\end{itemize}


\begin{itemize}
  \item \texttt{reactions}
  \item 渠道：Discord/Google Chat/Slack
  \item 必需：\texttt{--message-id}、\texttt{--target}
  \item 可选：\texttt{--limit}
\end{itemize}


\begin{itemize}
  \item \texttt{read}
  \item 渠道：Discord/Slack
  \item 必需：\texttt{--target}
  \item 可选：\texttt{--limit}、\texttt{--before}、\texttt{--after}
  \item 仅限 Discord：\texttt{--around}
\end{itemize}


\begin{itemize}
  \item \texttt{edit}
  \item 渠道：Discord/Slack
  \item 必需：\texttt{--message-id}、\texttt{--message}、\texttt{--target}
\end{itemize}


\begin{itemize}
  \item \texttt{delete}
  \item 渠道：Discord/Slack/Telegram
  \item 必需：\texttt{--message-id}、\texttt{--target}
\end{itemize}


\begin{itemize}
  \item \texttt{pin} / \texttt{unpin}
  \item 渠道：Discord/Slack
  \item 必需：\texttt{--message-id}、\texttt{--target}
\end{itemize}


\begin{itemize}
  \item \texttt{pins}（列表）
  \item 渠道：Discord/Slack
  \item 必需：\texttt{--target}
\end{itemize}


\begin{itemize}
  \item \texttt{permissions}
  \item 渠道：Discord
  \item 必需：\texttt{--target}
\end{itemize}


\begin{itemize}
  \item \texttt{search}
  \item 渠道：Discord
  \item 必需：\texttt{--guild-id}、\texttt{--query}
  \item 可选：\texttt{--channel-id}、\texttt{--channel-ids}（可重复）、\texttt{--author-id}、\texttt{--author-ids}（可重复）、\texttt{--limit}
\end{itemize}


\subsubsection{线程}


\begin{itemize}
  \item \texttt{thread create}
  \item 渠道：Discord
  \item 必需：\texttt{--thread-name}、\texttt{--target}（频道 ID）
  \item 可选：\texttt{--message-id}、\texttt{--auto-archive-min}
\end{itemize}


\begin{itemize}
  \item \texttt{thread list}
  \item 渠道：Discord
  \item 必需：\texttt{--guild-id}
  \item 可选：\texttt{--channel-id}、\texttt{--include-archived}、\texttt{--before}、\texttt{--limit}
\end{itemize}


\begin{itemize}
  \item \texttt{thread reply}
  \item 渠道：Discord
  \item 必需：\texttt{--target}（线程 ID）、\texttt{--message}
  \item 可选：\texttt{--media}、\texttt{--reply-to}
\end{itemize}


\subsubsection{表情符号}


\begin{itemize}
  \item \texttt{emoji list}
  \item Discord：\texttt{--guild-id}
  \item Slack：无需额外标志
\end{itemize}


\begin{itemize}
  \item \texttt{emoji upload}
  \item 渠道：Discord
  \item 必需：\texttt{--guild-id}、\texttt{--emoji-name}、\texttt{--media}
  \item 可选：\texttt{--role-ids}（可重复）
\end{itemize}


\subsubsection{贴纸}


\begin{itemize}
  \item \texttt{sticker send}
  \item 渠道：Discord
  \item 必需：\texttt{--target}、\texttt{--sticker-id}（可重复）
  \item 可选：\texttt{--message}
\end{itemize}


\begin{itemize}
  \item \texttt{sticker upload}
  \item 渠道：Discord
  \item 必需：\texttt{--guild-id}、\texttt{--sticker-name}、\texttt{--sticker-desc}、\texttt{--sticker-tags}、\texttt{--media}
\end{itemize}


\subsubsection{角色 / 频道 / 成员 / 语音}


\begin{itemize}
  \item \texttt{role info}（Discord）：\texttt{--guild-id}
  \item \texttt{role add} / \texttt{role remove}（Discord）：\texttt{--guild-id}、\texttt{--user-id}、\texttt{--role-id}
  \item \texttt{channel info}（Discord）：\texttt{--target}
  \item \texttt{channel list}（Discord）：\texttt{--guild-id}
  \item \texttt{member info}（Discord/Slack）：\texttt{--user-id}（Discord 还需要 \texttt{--guild-id}）
  \item \texttt{voice status}（Discord）：\texttt{--guild-id}、\texttt{--user-id}
\end{itemize}


\subsubsection{事件}


\begin{itemize}
  \item \texttt{event list}（Discord）：\texttt{--guild-id}
  \item \texttt{event create}（Discord）：\texttt{--guild-id}、\texttt{--event-name}、\texttt{--start-time}
  \item 可选：\texttt{--end-time}、\texttt{--desc}、\texttt{--channel-id}、\texttt{--location}、\texttt{--event-type}
\end{itemize}


\subsubsection{管理（Discord）}


\begin{itemize}
  \item \texttt{timeout}：\texttt{--guild-id}、\texttt{--user-id}（可选 \texttt{--duration-min} 或 \texttt{--until}；两者都省略则清除超时）
  \item \texttt{kick}：\texttt{--guild-id}、\texttt{--user-id}（+ \texttt{--reason}）
  \item \texttt{ban}：\texttt{--guild-id}、\texttt{--user-id}（+ \texttt{--delete-days}、\texttt{--reason}）
  \item \texttt{timeout} 也支持 \texttt{--reason}
\end{itemize}


\subsubsection{广播}


\begin{itemize}
  \item \texttt{broadcast}
  \item 渠道：任何已配置的渠道；使用 \texttt{--channel all} 可针对所有提供商
  \item 必需：\texttt{--targets}（可重复）
  \item 可选：\texttt{--message}、\texttt{--media}、\texttt{--dry-run}
\end{itemize}


\subsection{示例}


发送 Discord 回复：

\begin{lstlisting}
openclaw message send --channel discord \
  --target channel:123 --message "hi" --reply-to 456
\end{lstlisting}


创建 Discord 投票：

\begin{lstlisting}
openclaw message poll --channel discord \
  --target channel:123 \
  --poll-question "Snack?" \
  --poll-option Pizza --poll-option Sushi \
  --poll-multi --poll-duration-hours 48
\end{lstlisting}


发送 Teams 主动消息：

\begin{lstlisting}
openclaw message send --channel msteams \
  --target conversation:19:abc@thread.tacv2 --message "hi"
\end{lstlisting}


创建 Teams 投票：

\begin{lstlisting}
openclaw message poll --channel msteams \
  --target conversation:19:abc@thread.tacv2 \
  --poll-question "Lunch?" \
  --poll-option Pizza --poll-option Sushi
\end{lstlisting}


在 Slack 中添加表情回应：

\begin{lstlisting}
openclaw message react --channel slack \
  --target C123 --message-id 456 --emoji "✅"
\end{lstlisting}


在 Signal 群组中添加表情回应：

\begin{lstlisting}
openclaw message react --channel signal \
  --target signal:group:abc123 --message-id 1737630212345 \
  --emoji "✅" --target-author-uuid 123e4567-e89b-12d3-a456-426614174000
\end{lstlisting}


发送 Telegram 内联按钮：

\begin{lstlisting}
openclaw message send --channel telegram --target @mychat --message "Choose:" \
  --buttons '[ [{"text":"Yes","callback_data":"cmd:yes"}], [{"text":"No","callback_data":"cmd:no"}] ]'
\end{lstlisting}



\clearpage

% === 源文件: cli/models.md ===
\section{openclaw models}


模型发现、扫描和配置（默认模型、回退、认证配置）。

相关内容：

\begin{itemize}
  \item 提供商 + 模型：\href{/providers/models}{模型}
  \item 提供商认证设置：\href{/start/getting-started}{快速开始}
\end{itemize}


\subsection{常用命令}


\begin{lstlisting}[language=bash]
openclaw models status
openclaw models list
openclaw models set <model-or-alias>
openclaw models scan
\end{lstlisting}


\texttt{openclaw models status} 显示已解析的默认模型/回退配置以及认证概览。
当提供商使用快照可用时，OAuth/令牌状态部分会包含提供商使用头信息。
添加 \texttt{--probe} 可对每个已配置的提供商配置运行实时认证探测。
探测会发送真实请求（可能消耗令牌并触发速率限制）。
使用 \texttt{--agent } 可检查已配置智能体的模型/认证状态。省略时，
命令会使用 \texttt{OPENCLAW\_AGENT\_DIR}/\texttt{PI\_CODING\_AGENT\_DIR}（如已设置），否则使用
已配置的默认智能体。

注意事项：

\begin{itemize}
  \item \texttt{models set } 接受 \texttt{provider/model} 或别名。
  \item 模型引用通过在\textbf{第一个} \texttt{/} 处拆分来解析。如果模型 ID 包含 \texttt{/}（OpenRouter 风格），需包含提供商前缀（示例：\texttt{openrouter/moonshotai/kimi-k2}）。
  \item 如果省略提供商，OpenClaw 会将输入视为别名或\textbf{默认提供商}的模型（仅在模型 ID 不包含 \texttt{/} 时有效）。
\end{itemize}


\subsubsection{models status}


选项：

\begin{itemize}
  \item \texttt{--json}
  \item \texttt{--plain}
  \item \texttt{--check}（退出码 1=已过期/缺失，2=即将过期）
  \item \texttt{--probe}（对已配置的认证配置进行实时探测）
  \item \texttt{--probe-provider }（探测单个提供商）
  \item \texttt{--probe-profile }（可重复或逗号分隔的配置 ID）
  \item \texttt{--probe-timeout }
  \item \texttt{--probe-concurrency }
  \item \texttt{--probe-max-tokens }
  \item \texttt{--agent }（已配置的智能体 ID；覆盖 \texttt{OPENCLAW\_AGENT\_DIR}/\texttt{PI\_CODING\_AGENT\_DIR}）
\end{itemize}


\subsection{别名 + 回退}


\begin{lstlisting}[language=bash]
openclaw models aliases list
openclaw models fallbacks list
\end{lstlisting}


\subsection{认证配置}


\begin{lstlisting}[language=bash]
openclaw models auth add
openclaw models auth login --provider <id>
openclaw models auth setup-token
openclaw models auth paste-token
\end{lstlisting}


\texttt{models auth login} 运行提供商插件的认证流程（OAuth/API 密钥）。使用
\texttt{openclaw plugins list} 查看已安装的提供商。

注意事项：

\begin{itemize}
  \item \texttt{setup-token} 会提示输入 setup-token 值（在任意机器上使用 \texttt{claude setup-token} 生成）。
  \item \texttt{paste-token} 接受在其他地方或通过自动化生成的令牌字符串。
\end{itemize}



\clearpage

% === 源文件: cli/node.md ===
\section{openclaw node}


运行一个\textbf{无头节点主机}，连接到 Gateway 网关 WebSocket 并在此机器上暴露
\texttt{system.run} / \texttt{system.which}。

\subsection{为什么使用节点主机？}


当你希望智能体\textbf{在网络中的其他机器上运行命令}，而无需在那里安装完整的 macOS 配套应用时，请使用节点主机。

常见用例：

\begin{itemize}
  \item 在远程 Linux/Windows 机器上运行命令（构建服务器、实验室机器、NAS）。
  \item 在 Gateway 网关上保持执行的\textbf{沙箱隔离}，但将批准的运行委托给其他主机。
  \item 为自动化或 CI 节点提供轻量级、无头的执行目标。
\end{itemize}


执行仍然受\textbf{执行批准}和节点主机上的每智能体允许列表保护，因此你可以保持命令访问的范围明确。

\subsection{浏览器代理（零配置）}


如果节点上的 \texttt{browser.enabled} 未被禁用，节点主机会自动广播浏览器代理。这让智能体无需额外配置即可在该节点上使用浏览器自动化。

如需在节点上禁用：

\begin{lstlisting}[language=json]
{
  nodeHost: {
    browserProxy: {
      enabled: false,
    },
  },
}
\end{lstlisting}


\subsection{运行（前台）}


\begin{lstlisting}[language=bash]
openclaw node run --host <gateway-host> --port 18789
\end{lstlisting}


选项：

\begin{itemize}
  \item \texttt{--host }：Gateway 网关 WebSocket 主机（默认：\texttt{127.0.0.1}）
  \item \texttt{--port }：Gateway 网关 WebSocket 端口（默认：\texttt{18789}）
  \item \texttt{--tls}：为 Gateway 网关连接使用 TLS
  \item \texttt{--tls-fingerprint }：预期的 TLS 证书指纹（sha256）
  \item \texttt{--node-id }：覆盖节点 id（清除配对 token）
  \item \texttt{--display-name }：覆盖节点显示名称
\end{itemize}


\subsection{服务（后台）}


将无头节点主机安装为用户服务。

\begin{lstlisting}[language=bash]
openclaw node install --host <gateway-host> --port 18789
\end{lstlisting}


选项：

\begin{itemize}
  \item \texttt{--host }：Gateway 网关 WebSocket 主机（默认：\texttt{127.0.0.1}）
  \item \texttt{--port }：Gateway 网关 WebSocket 端口（默认：\texttt{18789}）
  \item \texttt{--tls}：为 Gateway 网关连接使用 TLS
  \item \texttt{--tls-fingerprint }：预期的 TLS 证书指纹（sha256）
  \item \texttt{--node-id }：覆盖节点 id（清除配对 token）
  \item \texttt{--display-name }：覆盖节点显示名称
  \item \texttt{--runtime }：服务运行时（\texttt{node} 或 \texttt{bun}）
  \item \texttt{--force}：如果已安装则重新安装/覆盖
\end{itemize}


管理服务：

\begin{lstlisting}[language=bash]
openclaw node status
openclaw node stop
openclaw node restart
openclaw node uninstall
\end{lstlisting}


使用 \texttt{openclaw node run} 运行前台节点主机（无服务）。

服务命令接受 \texttt{--json} 以获取机器可读输出。

\subsection{配对}


首次连接会在 Gateway 网关上创建待处理的节点配对请求。
通过以下方式批准：

\begin{lstlisting}[language=bash]
openclaw nodes pending
openclaw nodes approve <requestId>
\end{lstlisting}


节点主机将其节点 id、token、显示名称和 Gateway 网关连接信息存储在
\texttt{\textasciitilde{}/.openclaw/node.json} 中。

\subsection{执行批准}


\texttt{system.run} 受本地执行批准限制：

\begin{itemize}
  \item \texttt{\textasciitilde{}/.openclaw/exec-approvals.json}
  \item \href{/tools/exec-approvals}{执行批准}
  \item \texttt{openclaw approvals --node }（从 Gateway 网关编辑）
\end{itemize}



\clearpage

% === 源文件: cli/nodes.md ===
\section{openclaw nodes}


管理已配对的节点（设备）并调用节点功能。

相关内容：

\begin{itemize}
  \item 节点概述：\href{/nodes}{节点}
  \item 摄像头：\href{/nodes/camera}{摄像头节点}
  \item 图像：\href{/nodes/images}{图像节点}
\end{itemize}


通用选项：

\begin{itemize}
  \item \texttt{--url}、\texttt{--token}、\texttt{--timeout}、\texttt{--json}
\end{itemize}


\subsection{常用命令}


\begin{lstlisting}[language=bash]
openclaw nodes list
openclaw nodes list --connected
openclaw nodes list --last-connected 24h
openclaw nodes pending
openclaw nodes approve <requestId>
openclaw nodes status
openclaw nodes status --connected
openclaw nodes status --last-connected 24h
\end{lstlisting}


\texttt{nodes list} 打印待处理/已配对表格。已配对行包含最近连接时长（Last Connect）。
使用 \texttt{--connected} 仅显示当前已连接的节点。使用 \texttt{--last-connected }
筛选在指定时间段内连接过的节点（例如 \texttt{24h}、\texttt{7d}）。

\subsection{调用 / 运行}


\begin{lstlisting}[language=bash]
openclaw nodes invoke --node <id|name|ip> --command <command> --params <json>
openclaw nodes run --node <id|name|ip> <command...>
openclaw nodes run --raw "git status"
openclaw nodes run --agent main --node <id|name|ip> --raw "git status"
\end{lstlisting}


调用标志：

\begin{itemize}
  \item \texttt{--params }：JSON 对象字符串（默认 \texttt{\{\}}）。
  \item \texttt{--invoke-timeout }：节点调用超时（默认 \texttt{15000}）。
  \item \texttt{--idempotency-key }：可选的幂等键。
\end{itemize}


\subsubsection{Exec 风格默认值}


\texttt{nodes run} 与模型的 exec 行为一致（默认值 + 审批）：

\begin{itemize}
  \item 读取 \texttt{tools.exec.\textit{}（以及 \texttt{agents.list[].tools.exec.}} 覆盖）。
  \item 在调用 \texttt{system.run} 前使用 exec 审批（\texttt{exec.approval.request}）。
  \item 当设置了 \texttt{tools.exec.node} 时可省略 \texttt{--node}。
  \item 需要支持 \texttt{system.run} 的节点（macOS 配套应用或无头节点主机）。
\end{itemize}


标志：

\begin{itemize}
  \item \texttt{--cwd }：工作目录。
  \item \texttt{--env }：环境变量覆盖（可重复）。
  \item \texttt{--command-timeout }：命令超时。
  \item \texttt{--invoke-timeout }：节点调用超时（默认 \texttt{30000}）。
  \item \texttt{--needs-screen-recording}：要求屏幕录制权限。
  \item \texttt{--raw }：运行 shell 字符串（\texttt{/bin/sh -lc} 或 \texttt{cmd.exe /c}）。
  \item \texttt{--agent }：智能体范围的审批/白名单（默认为已配置的智能体）。
  \item \texttt{--ask }、\texttt{--security }：覆盖选项。
\end{itemize}



\clearpage

% === 源文件: cli/onboard.md ===
\section{openclaw onboard}


交互式新手引导向导（本地或远程 Gateway 网关设置）。

相关内容：

\begin{itemize}
  \item 向导指南：\href{/start/onboarding}{新手引导}
\end{itemize}


\subsection{示例}


\begin{lstlisting}[language=bash]
openclaw onboard
openclaw onboard --flow quickstart
openclaw onboard --flow manual
openclaw onboard --mode remote --remote-url ws://gateway-host:18789
\end{lstlisting}


流程说明：

\begin{itemize}
  \item \texttt{quickstart}：最少提示，自动生成 Gateway 网关令牌。
  \item \texttt{manual}：完整的端口/绑定/认证提示（\texttt{advanced} 的别名）。
  \item 最快开始聊天：\texttt{openclaw dashboard}（控制 UI，无需渠道设置）。
\end{itemize}



\clearpage

% === 源文件: cli/pairing.md ===
\section{openclaw pairing}


批准或检查私信配对请求（适用于支持配对的渠道）。

相关内容：

\begin{itemize}
  \item 配对流程：\href{/channels/pairing}{配对}
\end{itemize}


\subsection{命令}


\begin{lstlisting}[language=bash]
openclaw pairing list whatsapp
openclaw pairing approve whatsapp <code> --notify
\end{lstlisting}



\clearpage

% === 源文件: cli/plugins.md ===
\section{openclaw plugins}


管理 Gateway 网关插件/扩展（进程内加载）。

相关内容：

\begin{itemize}
  \item 插件系统：\href{/tools/plugin}{插件}
  \item 插件清单 + 模式：\href{/plugins/manifest}{插件清单}
  \item 安全加固：\href{/gateway/security}{安全}
\end{itemize}


\subsection{命令}


\begin{lstlisting}[language=bash]
openclaw plugins list
openclaw plugins info <id>
openclaw plugins enable <id>
openclaw plugins disable <id>
openclaw plugins doctor
openclaw plugins update <id>
openclaw plugins update --all
\end{lstlisting}


内置插件随 OpenClaw 一起发布，但默认禁用。使用 \texttt{plugins enable} 来激活它们。

所有插件必须提供 \texttt{openclaw.plugin.json} 文件，其中包含内联 JSON Schema（\texttt{configSchema}，即使为空）。缺少或无效的清单或模式会阻止插件加载并导致配置验证失败。

\subsubsection{安装}


\begin{lstlisting}[language=bash]
openclaw plugins install <path-or-spec>
\end{lstlisting}


安全提示：将插件安装视为运行代码。优先使用固定版本。

支持的归档格式：\texttt{.zip}、\texttt{.tgz}、\texttt{.tar.gz}、\texttt{.tar}。

使用 \texttt{--link} 避免复制本地目录（添加到 \texttt{plugins.load.paths}）：

\begin{lstlisting}[language=bash]
openclaw plugins install -l ./my-plugin
\end{lstlisting}


\subsubsection{更新}


\begin{lstlisting}[language=bash]
openclaw plugins update <id>
openclaw plugins update --all
openclaw plugins update <id> --dry-run
\end{lstlisting}


更新仅适用于从 npm 安装的插件（在 \texttt{plugins.installs} 中跟踪）。


\clearpage

% === 源文件: cli/reset.md ===
\section{openclaw reset}


重置本地配置/状态（保留 CLI 安装）。

\begin{lstlisting}[language=bash]
openclaw reset
openclaw reset --dry-run
openclaw reset --scope config+creds+sessions --yes --non-interactive
\end{lstlisting}



\clearpage

% === 源文件: cli/sandbox.md ===
\section{沙箱 CLI}


管理基于 Docker 的沙箱容器，用于隔离智能体执行。

\subsection{概述}


OpenClaw 可以在隔离的 Docker 容器中运行智能体以确保安全。\texttt{sandbox} 命令帮助你管理这些容器，特别是在更新或配置更改后。

\subsection{命令}


\subsubsection{openclaw sandbox explain}


检查\textbf{生效的}沙箱模式/作用域/工作区访问权限、沙箱工具策略和提权门控（附带修复配置的键路径）。

\begin{lstlisting}[language=bash]
openclaw sandbox explain
openclaw sandbox explain --session agent:main:main
openclaw sandbox explain --agent work
openclaw sandbox explain --json
\end{lstlisting}


\subsubsection{openclaw sandbox list}


列出所有沙箱容器及其状态和配置。

\begin{lstlisting}[language=bash]
openclaw sandbox list
openclaw sandbox list --browser  # List only browser containers
openclaw sandbox list --json     # JSON output
\end{lstlisting}


\textbf{输出包括：}

\begin{itemize}
  \item 容器名称和状态（运行中/已停止）
  \item Docker 镜像及其是否与配置匹配
  \item 创建时间
  \item 空闲时间（自上次使用以来的时间）
  \item 关联的会话/智能体
\end{itemize}


\subsubsection{openclaw sandbox recreate}


移除沙箱容器以强制使用更新的镜像/配置重新创建。

\begin{lstlisting}[language=bash]
openclaw sandbox recreate --all                # Recreate all containers
openclaw sandbox recreate --session main       # Specific session
openclaw sandbox recreate --agent mybot        # Specific agent
openclaw sandbox recreate --browser            # Only browser containers
openclaw sandbox recreate --all --force        # Skip confirmation
\end{lstlisting}


\textbf{选项：}

\begin{itemize}
  \item \texttt{--all}：重新创建所有沙箱容器
  \item \texttt{--session }：重新创建特定会话的容器
  \item \texttt{--agent }：重新创建特定智能体的容器
  \item \texttt{--browser}：仅重新创建浏览器容器
  \item \texttt{--force}：跳过确认提示
\end{itemize}


\textbf{重要：} 容器会在智能体下次使用时自动重新创建。

\subsection{使用场景}


\subsubsection{更新 Docker 镜像后}


\begin{lstlisting}[language=bash]
# Pull new image
docker pull openclaw-sandbox:latest
docker tag openclaw-sandbox:latest openclaw-sandbox:bookworm-slim

# Update config to use new image
# Edit config: agents.defaults.sandbox.docker.image (or agents.list[].sandbox.docker.image)

# Recreate containers
openclaw sandbox recreate --all
\end{lstlisting}


\subsubsection{更改沙箱配置后}


\begin{lstlisting}[language=bash]
# Edit config: agents.defaults.sandbox.* (or agents.list[].sandbox.*)

# Recreate to apply new config
openclaw sandbox recreate --all
\end{lstlisting}


\subsubsection{更改 setupCommand 后}


\begin{lstlisting}[language=bash]
openclaw sandbox recreate --all
# or just one agent:
openclaw sandbox recreate --agent family
\end{lstlisting}


\subsubsection{仅针对特定智能体}


\begin{lstlisting}[language=bash]
# Update only one agent's containers
openclaw sandbox recreate --agent alfred
\end{lstlisting}


\subsection{为什么需要这个？}


\textbf{问题：} 当你更新沙箱 Docker 镜像或配置时：

\begin{itemize}
  \item 现有容器继续使用旧设置运行
  \item 容器仅在空闲 24 小时后才被清理
  \item 经常使用的智能体会无限期保持旧容器运行
\end{itemize}


\textbf{解决方案：} 使用 \texttt{openclaw sandbox recreate} 强制移除旧容器。它们会在下次需要时自动使用当前设置重新创建。

提示：优先使用 \texttt{openclaw sandbox recreate} 而不是手动 \texttt{docker rm}。它使用 Gateway 网关的容器命名规则，避免在作用域/会话键更改时出现不匹配。

\subsection{配置}


沙箱设置位于 \texttt{\textasciitilde{}/.openclaw/openclaw.json} 的 \texttt{agents.defaults.sandbox} 下（每个智能体的覆盖设置在 \texttt{agents.list[].sandbox} 中）：

\begin{lstlisting}[language=json]
{
  "agents": {
    "defaults": {
      "sandbox": {
        "mode": "all", // off, non-main, all
        "scope": "agent", // session, agent, shared
        "docker": {
          "image": "openclaw-sandbox:bookworm-slim",
          "containerPrefix": "openclaw-sbx-",
          // ... more Docker options
        },
        "prune": {
          "idleHours": 24, // Auto-prune after 24h idle
          "maxAgeDays": 7, // Auto-prune after 7 days
        },
      },
    },
  },
}
\end{lstlisting}


\subsection{另请参阅}


\begin{itemize}
  \item \href{/gateway/sandboxing}{沙箱文档}
  \item \href{/concepts/agent-workspace}{智能体配置}
  \item \href{/gateway/doctor}{Doctor 命令} - 检查沙箱设置
\end{itemize}



\clearpage

% === 源文件: cli/security.md ===
\section{openclaw security}


安全工具（审计 + 可选修复）。

相关：

\begin{itemize}
  \item 安全指南：\href{/gateway/security}{安全}
\end{itemize}


\subsection{审计}


\begin{lstlisting}[language=bash]
openclaw security audit
openclaw security audit --deep
openclaw security audit --fix
\end{lstlisting}


当多个私信发送者共享主会话时，审计会发出警告，并建议对共享收件箱使用 \texttt{session.dmScope="per-channel-peer"}（或多账户渠道使用 \texttt{per-account-channel-peer}）。
当使用小模型（\texttt{<=300B}）且未启用沙箱隔离但启用了 web/browser 工具时，它也会发出警告。


\clearpage

% === 源文件: cli/sessions.md ===
\section{openclaw sessions}


列出已存储的对话会话。

\begin{lstlisting}[language=bash]
openclaw sessions
openclaw sessions --active 120
openclaw sessions --json
\end{lstlisting}



\clearpage

% === 源文件: cli/setup.md ===
\section{openclaw setup}


初始化 \texttt{\textasciitilde{}/.openclaw/openclaw.json} 和智能体工作区。

相关内容：

\begin{itemize}
  \item 快速开始：\href{/start/getting-started}{快速开始}
  \item 向导：\href{/start/onboarding}{新手引导}
\end{itemize}


\subsection{示例}


\begin{lstlisting}[language=bash]
openclaw setup
openclaw setup --workspace ~/.openclaw/workspace
\end{lstlisting}


通过 setup 运行向导：

\begin{lstlisting}[language=bash]
openclaw setup --wizard
\end{lstlisting}



\clearpage

% === 源文件: cli/skills.md ===
\section{openclaw skills}


检查 Skills（内置 + 工作区 + 托管覆盖）并查看哪些符合条件，哪些缺少要求。

相关内容：

\begin{itemize}
  \item Skills 系统：\href{/tools/skills}{Skills}
  \item Skills 配置：\href{/tools/skills-config}{Skills 配置}
  \item ClawHub 安装：\href{/tools/clawhub}{ClawHub}
\end{itemize}


\subsection{命令}


\begin{lstlisting}[language=bash]
openclaw skills list
openclaw skills list --eligible
openclaw skills info <name>
openclaw skills check
\end{lstlisting}



\clearpage

% === 源文件: cli/status.md ===
\section{openclaw status}


渠道 + 会话的诊断。

\begin{lstlisting}[language=bash]
openclaw status
openclaw status --all
openclaw status --deep
openclaw status --usage
\end{lstlisting}


注意事项：

\begin{itemize}
  \item \texttt{--deep} 运行实时探测（WhatsApp Web + Telegram + Discord + Google Chat + Slack + Signal）。
  \item 当配置了多个智能体时，输出包含每个智能体的会话存储。
  \item 概览包含 Gateway 网关 + 节点主机服务安装/运行时状态（如果可用）。
  \item 概览包含更新渠道 + git SHA（用于源代码检出）。
  \item 更新信息显示在概览中；如果有可用更新，status 会打印提示运行 \texttt{openclaw update}（参见\href{/install/updating}{更新}）。
\end{itemize}



\clearpage

% === 源文件: cli/system.md ===
\section{openclaw system}


Gateway 网关的系统级辅助工具：入队系统事件、控制心跳和查看在线状态。

\subsection{常用命令}


\begin{lstlisting}[language=bash]
openclaw system event --text "Check for urgent follow-ups" --mode now
openclaw system heartbeat enable
openclaw system heartbeat last
openclaw system presence
\end{lstlisting}


\subsection{system event}


在\textbf{主}会话上入队系统事件。下一次心跳会将其作为 \texttt{System:} 行注入到提示中。使用 \texttt{--mode now} 立即触发心跳；\texttt{next-heartbeat} 等待下一个计划的心跳时刻。

标志：

\begin{itemize}
  \item \texttt{--text }：必填的系统事件文本。
  \item \texttt{--mode }：\texttt{now} 或 \texttt{next-heartbeat}（默认）。
  \item \texttt{--json}：机器可读输出。
\end{itemize}


\subsection{system heartbeat last|enable|disable}


心跳控制：

\begin{itemize}
  \item \texttt{last}：显示最后一次心跳事件。
  \item \texttt{enable}：重新开启心跳（如果之前被禁用，使用此命令）。
  \item \texttt{disable}：暂停心跳。
\end{itemize}


标志：

\begin{itemize}
  \item \texttt{--json}：机器可读输出。
\end{itemize}


\subsection{system presence}


列出 Gateway 网关已知的当前系统在线状态条目（节点、实例和类似状态行）。

标志：

\begin{itemize}
  \item \texttt{--json}：机器可读输出。
\end{itemize}


\subsection{注意}


\begin{itemize}
  \item 需要一个运行中的 Gateway 网关，可通过你当前的配置访问（本地或远程）。
  \item 系统事件是临时的，不会在重启后持久化。
\end{itemize}



\clearpage

% === 源文件: cli/tui.md ===
\section{openclaw tui}


打开连接到 Gateway 网关的终端 UI。

相关：

\begin{itemize}
  \item TUI 指南：\href{/web/tui}{TUI}
\end{itemize}


\subsection{示例}


\begin{lstlisting}[language=bash]
openclaw tui
openclaw tui --url ws://127.0.0.1:18789 --token <token>
openclaw tui --session main --deliver
\end{lstlisting}



\clearpage

% === 源文件: cli/uninstall.md ===
\section{openclaw uninstall}


卸载 Gateway 网关服务 + 本地数据（CLI 保留）。

\begin{lstlisting}[language=bash]
openclaw uninstall
openclaw uninstall --all --yes
openclaw uninstall --dry-run
\end{lstlisting}



\clearpage

% === 源文件: cli/update.md ===
\section{openclaw update}


安全更新 OpenClaw 并在 stable/beta/dev 渠道之间切换。

如果你通过 \textbf{npm/pnpm} 安装（全局安装，无 git 元数据），更新通过 \href{/install/updating}{更新} 中的包管理器流程进行。

\subsection{用法}


\begin{lstlisting}[language=bash]
openclaw update
openclaw update status
openclaw update wizard
openclaw update --channel beta
openclaw update --channel dev
openclaw update --tag beta
openclaw update --no-restart
openclaw update --json
openclaw --update
\end{lstlisting}


\subsection{选项}


\begin{itemize}
  \item \texttt{--no-restart}：成功更新后跳过重启 Gateway 网关服务。
  \item \texttt{--channel }：设置更新渠道（git + npm；持久化到配置中）。
  \item \texttt{--tag }：仅为本次更新覆盖 npm dist-tag 或版本。
  \item \texttt{--json}：打印机器可读的 \texttt{UpdateRunResult} JSON。
  \item \texttt{--timeout }：每步超时时间（默认 1200 秒）。
\end{itemize}


注意：降级需要确认，因为旧版本可能会破坏配置。

\subsection{update status}


显示当前更新渠道 + git 标签/分支/SHA（对于源码检出），以及更新可用性。

\begin{lstlisting}[language=bash]
openclaw update status
openclaw update status --json
openclaw update status --timeout 10
\end{lstlisting}


选项：

\begin{itemize}
  \item \texttt{--json}：打印机器可读的状态 JSON。
  \item \texttt{--timeout }：检查超时时间（默认 3 秒）。
\end{itemize}


\subsection{update wizard}


交互式流程，用于选择更新渠道并确认是否在更新后重启 Gateway 网关（默认重启）。如果你选择 \texttt{dev} 但没有 git 检出，它会提供创建一个的选项。

\subsection{工作原理}


当你显式切换渠道（\texttt{--channel ...}）时，OpenClaw 也会保持安装方式一致：

\begin{itemize}
  \item \texttt{dev} → 确保存在 git 检出（默认：\texttt{\textasciitilde{}/openclaw}，可通过 \texttt{OPENCLAW\_GIT\_DIR} 覆盖），更新它，并从该检出安装全局 CLI。
  \item \texttt{stable}/\texttt{beta} → 使用匹配的 dist-tag 从 npm 安装。
\end{itemize}


\subsection{Git 检出流程}


渠道：

\begin{itemize}
  \item \texttt{stable}：检出最新的非 beta 标签，然后构建 + doctor。
  \item \texttt{beta}：检出最新的 \texttt{-beta} 标签，然后构建 + doctor。
  \item \texttt{dev}：检出 \texttt{main}，然后 fetch + rebase。
\end{itemize}


高层概述：

\begin{enumerate}
  \item 需要干净的工作树（无未提交的更改）。
  \item 切换到所选渠道（标签或分支）。
  \item 获取上游（仅 dev）。
  \item 仅 dev：在临时工作树中预检 lint + TypeScript 构建；如果最新提交失败，回退最多 10 个提交以找到最新的干净构建。
  \item Rebase 到所选提交（仅 dev）。
  \item 安装依赖（优先使用 pnpm；npm 作为备选）。
  \item 构建 + 构建控制界面。
  \item 运行 \texttt{openclaw doctor} 作为最终的"安全更新"检查。
  \item 将插件同步到当前渠道（dev 使用捆绑的扩展；stable/beta 使用 npm）并更新 npm 安装的插件。
\end{enumerate}


\subsection{--update 简写}


\texttt{openclaw --update} 会重写为 \texttt{openclaw update}（便于 shell 和启动脚本使用）。

\subsection{另请参阅}


\begin{itemize}
  \item \texttt{openclaw doctor}（在 git 检出上会提供先运行更新的选项）
  \item \href{/install/development-channels}{开发渠道}
  \item \href{/install/updating}{更新}
  \item \href{/cli}{CLI 参考}
\end{itemize}



\clearpage

% === 源文件: cli/voicecall.md ===
\section{openclaw voicecall}


\texttt{voicecall} 是一个由插件提供的命令。只有在安装并启用了语音通话插件时才会出现。

主要文档：

\begin{itemize}
  \item 语音通话插件：\href{/plugins/voice-call}{语音通话}
\end{itemize}


\subsection{常用命令}


\begin{lstlisting}[language=bash]
openclaw voicecall status --call-id <id>
openclaw voicecall call --to "+15555550123" --message "Hello" --mode notify
openclaw voicecall continue --call-id <id> --message "Any questions?"
openclaw voicecall end --call-id <id>
\end{lstlisting}


\subsection{暴露 Webhook（Tailscale）}


\begin{lstlisting}[language=bash]
openclaw voicecall expose --mode serve
openclaw voicecall expose --mode funnel
openclaw voicecall unexpose
\end{lstlisting}


安全提示：仅将 webhook 端点暴露给你信任的网络。尽可能优先使用 Tailscale Serve 而非 Funnel。


\clearpage

% === 源文件: cli/webhooks.md ===
\section{openclaw webhooks}


Webhook 辅助工具和集成（Gmail Pub/Sub、Webhook 辅助工具）。

相关内容：

\begin{itemize}
  \item Webhook：\href{/automation/webhook}{Webhook}
  \item Gmail Pub/Sub：\href{/automation/gmail-pubsub}{Gmail Pub/Sub}
\end{itemize}


\subsection{Gmail}


\begin{lstlisting}[language=bash]
openclaw webhooks gmail setup --account you@example.com
openclaw webhooks gmail run
\end{lstlisting}


详情请参阅 \href{/automation/gmail-pubsub}{Gmail Pub/Sub 文档}。


\clearpage
